% !TEX root = print.tex
% The blueprint's content, split by chapter. Each part holds the statement nodes; this
% file only orders them. See the framework's LINKAGE.md for the node conventions.
%
% ---------------------------------------------------------------------------------
% NUMBERING: one chapter per article section, chapter number == the article's section
% number. REORDERED 2026-09-01 (publication plan P2, author's decision).
% ---------------------------------------------------------------------------------
%
% theorems.tex masters one shared counter off `definition`, reset per chapter. The
% chapters below follow the ARTICLE's order (paper/main.tex): axioms, preliminaries,
% representation, cascade, covariance, characterization, examples, signaling, locality,
% implementation, jet, with the memory-kernel/scale-Cauchy pair as the appendix it is in
% the article. Starting the chapter counter at 1 makes the blueprint's chapter numbers
% match the paper's section numbers -- the signaling chapter is Chapter 9 as the
% signaling section is \S9 -- so the two public documents read side by side.
%
% The original numbering was draft-aligned (one chapter per draft section, so that
% Lemma 4.1 was the draft's Lemma 4.1). That goal is RETIRED (author, 2026-09-01): the
% article is the deliverable, and draft alignment served the transcription, which is
% done. draft/ stays in the repo as the starting point of record, and every node still
% carries a `% draft: <kind> N.M` provenance comment -- that comment, not the rendered
% counter, is now the bridge to the draft. Historical comments in the parts that cite
% node numbers ("Numbered 11.25") or chapter numbers cite the retired draft-aligned
% numbering; each part's header says what it renders as now. The part FILES keep their
% draft-numbered names (11-signaling.tex renders as Chapter 9): the names are provenance,
% and renaming them would break the draft bridge for no rendered gain.
%
% There is still no Chapter 1: \S1 of the article is the introduction, positioning
% rather than mathematics, and is not transcribed. Chapter 2 is the axiom system, as \S2
% of the article is temporal measurement (the paper's motivating prose stays there).
%
% Scope is the article's mathematics, \S\S2-12 plus the appendix. \S13 of the article
% (concluding remarks) is not transcribed. Scope note of record (2026-09-01, author's
% instruction): the implementation and jet chapters were originally excluded to keep the
% blueprint's size under control; the author wants the blueprint complete with respect to
% the mathematical material, with formalisation a separate decision -- so those chapters
% carry statement nodes without \lean tags, and none is scheduled in a Lean phase.

% The .tex extensions below are explicit for the reader's sake; all three consumers
% accept either form (pdflatex natively, `linkage` by appending .tex when absent,
% plasTeX through kpsewhich).
%
% What IS load-bearing, and cost several CI runs to find: plasTeX resolves every
% \input through kpsewhich, and its fallback when that binary is missing is an exact
% filename match against a directory listing -- which can never resolve a name
% containing a separator. So on a machine without a working kpsewhich, web.tex and
% content.tex resolve while every parts/* file is silently dropped, and the web build
% renders a complete-looking site with no statement nodes in it. The shared docs
% workflow installs texlive-base for this and asserts kpsewhich before building.
% Adding .tex does not help; the separator is the issue.
% !TEX root = print.tex
% This article's additional theorem-like environments.
%
% Framework-owned theorems.tex defines definition/lemma/proposition/theorem on one
% counter mastered by `definition` and reset per chapter. Three more are needed here,
% and all three ride the SAME counter, because the numbering scheme below depends on it.
%
% Why they are not in macros.tex: both drivers (print.tex, web.tex) \input macros
% *before* theorems, so `definition` does not exist yet at that point. This file is
% \input from the top of content.tex instead, after theorems.tex has run.
%
%   corollary -- omitted by theorems.tex although `linkage` counts it a statement
%                environment (config.py DEFAULT_STATEMENT_ENVS) and plastexdepgraph
%                graphs it by default. The draft has three. A `corollary` here IS a
%                blueprint node and needs \label/\statusT/\statusA like any other.
%
%   remark    -- NOT a statement environment: no \label discipline, no status tag, not
%   example      in the dependency graph. They carry the draft's Remarks and Examples so
%                the blueprint stays the full text of record. A remark or example whose
%                content is load-bearing must be promoted to a definition/proposition
%                instead -- see the \S8 families, where the draft's Examples 8.1-8.3
%                become nodes.
%
% Adding a *statement* kind beyond corollary would also mean listing it in web.tex's
% `thms=` option, or the graph would silently ignore its \uses/\leanok annotations.

\newtheorem{corollary}[definition]{Corollary}
\newtheorem{remark}[definition]{Remark}
\newtheorem{example}[definition]{Example}


\setcounter{chapter}{1}

\chapter{Axioms}
% !TEX root = ../print.tex
% Draft \S3 --- Axioms
%
% Source: draft/hemigroup-causal-scale-space-kernels.md, lines 68-92.
% This file is chapter 3 --- Axioms of the blueprint; the shared counter makes every number here
% agree with the draft's own (see content.tex).
%
% Planned nodes: def:cascade-family (A1)-(A8)+(ND); Remark 3.2
%
% No nodes yet. Each statement carries \label{kind:name}, \lean{...} with
% \leanok/\notready, \uses{...}, \notes{slug}, a \statusT or \statusA line, and a
% "% draft: <kind> N.M" provenance comment; an \statusA node must also declare what its
% citation carries (\textbf{Assignment.}).


\chapter{Preliminaries and notation}
% !TEX root = ../print.tex
% Draft \S2 --- Preliminaries and notation
%
% Source: draft/hemigroup-causal-scale-space-kernels.md, lines 31-67.
% Chapter 2 of the blueprint; the shared counter puts every number here on the draft's own.
%
% TWO DEPARTURES FROM THE DRAFT, both deliberate:
%
% 1. Proposition 2.3 stays ONE node with four parts, rather than four nodes. Splitting it
%    would push Lemma 2.4 to 2.7 and break the numbering alignment for the section the rest
%    of the article cites most. The per-clause honesty lives in the \textbf{Assignment.}
%    clause instead, which is what that mechanism is for (LINKAGE.md rule 7).
%
% 2. prop:laplace-continuity (2.5) and prop:laplace-uniqueness (2.6) are ADDITIVE -- the
%    draft uses both facts (Lemma 4.1's uniqueness clause, Lemma 6.1, Theorem 7.3's (A7)
%    verification, Proposition 9.2, Theorem 9.5, Lemma 10.3(5)) but never states them, so
%    their trust was invisible to the dependency graph. Appending them leaves 2.1-2.4 alone.
%
% 3. 2.7-2.10 are ADDITIVE for the same reason and by the same rule -- appended, so that
%    2.1-2.6 keep the draft's numbers. They exist because the FORMALISATION had outrun the
%    graph: ~900 lines of finished, axiom-guarded Lean (Levy, Injectivity, WeakConvergence,
%    the Markov bound of Continuity) proved things this chapter either did not state at all
%    (the representation, Def. 2.7 -- the single most-used object in the Lean development,
%    and until now nameable only inside Prop. 2.3(3)) or stated only in a stronger cited
%    form it could not be tagged against (2.5, 2.6). The pattern for the latter is worth
%    naming, because chapters 8-12 will want it again:
%
%      an [A] node keeps the PAPER's statement and its ledger entry; a [T] node beneath it
%      states the RESTRICTED case the development actually proves, carries the \lean tag,
%      and is what the downstream node \uses.
%
%    So Prop. 2.5 and 2.6 are unchanged -- the paper still cites Feller -- while 2.8 and 2.9
%    record that nothing formalised here rests on A5 or A6.
% \notes{} tags are omitted throughout until the hub's content-note slugs exist; an
% unresolvable slug fails `linkage check --wiki`.

Time runs over $\RR$. We write $X = \Lone$, real, with positive cone $X_+$ and
$\int f := \int_{\RR} f(t)\,dt$; $\Mplus(\RR)$ for the finite positive Borel measures on
$\RR$, with total mass $\|\mu\| = \mu(\RR)$; translations $(\shift_a f)(t) = f(t-a)$ for
$a \in \RR$; dilations on functions $(\dil_\sigma f)(t) = \sigma^{-1} f(t/\sigma)$,
$\sigma > 0$, which preserve mass, and on measures $\dil_\sigma \mu :=$ the pushforward of
$\mu$ under $t \mapsto \sigma t$; and convolution
$(\mu * f)(t) = \int f(t - r)\,\mu(dr)$ for $\mu \in \Mplus(\RR)$, $f \in X$. Then
$\|\mu * f\|_1 \le \|\mu\|\,\|f\|_1$, and dilation is compatible with convolution:
\[
  \dil_\sigma(\mu * f) = (\dil_\sigma \mu) * (\dil_\sigma f).
\]

For $\mu \in \Mplus$ supported on $[0,\infty)$ the Laplace transform is
\[
  \hat\mu(s) = \int_{[0,\infty)} e^{-st}\,\mu(dt), \qquad s \ge 0 .
\]
If $\mu \ne 0$ then $\hat\mu(s) > 0$ for all $s \ge 0$.

% draft: Definition 2.1
\begin{definition}[completely monotone]\label{def:completely-monotone}
A function $f : (0,\infty) \to \RR$ is \emph{completely monotone} (CM) if $f \in C^\infty$
and $(-1)^n f^{(n)} \ge 0$ for all $n \ge 0$.
\statusT
\end{definition}

% draft: Definition 2.2
\begin{definition}[Bernstein function]\label{def:bernstein-function}
\uses{def:completely-monotone}
A function $g : (0,\infty) \to [0,\infty)$ is a \emph{Bernstein function} if $g \in C^\infty$
and $g'$ is completely monotone. We write $g \in \BF$, and $g \in \BFz$ if moreover
$g(\zp) = 0$.
\statusT
\end{definition}

% draft: Proposition 2.3 (toolbox)
\begin{proposition}[the Bernstein-function toolbox]\label{prop:bernstein-toolbox}
\uses{def:completely-monotone,def:bernstein-function}
\begin{enumerate}
\item (Bernstein--Widder / Feller criterion.) $f$ is the Laplace transform of a probability
measure on $[0,\infty)$ if and only if $f$ is completely monotone and $f(\zp) = 1$. More
generally, $f$ is completely monotone if and only if $f(s) = \int_{[0,\infty)} e^{-st}\,F(dt)$
for a measure $F$ on $[0,\infty)$, not necessarily finite, and $F$ is then unique.
\item For $g : (0,\infty) \to [0,\infty)$ the following are equivalent: $g \in \BF$; the
composition $h \circ g$ is completely monotone for every completely monotone $h$; and
$e^{-\tau g}$ is completely monotone for every $\tau > 0$.
\item Every $g \in \BF$ is nonnegative, nondecreasing and concave, and has the
Lévy--Khintchine representation
\[
  g(s) = a + b s + \int_0^\infty \bigl(1 - e^{-st}\bigr)\,\nu(dt), \qquad
  a, b \ge 0, \quad \int_0^\infty (1 \wedge t)\,\nu(dt) < \infty,
\]
with $g(\zp) = a$. The triple $(a, b, \nu)$ is unique.
\item $\BF$ is a convex cone, closed under pointwise limits: if $g_n \in \BF$ and
$g_n(s) \to g(s) < \infty$ for every $s > 0$, then $g \in \BF$. Consequently $\BF$ is closed
under mixtures $s \mapsto \int g_u(s)\,m(du)$ with $m \ge 0$, whenever the integral is finite.
\end{enumerate}
\statusA\quad\ledger{A1}\quad\ledger{A2}\quad\ledger{A3}\quad\ledger{A4}\quad\emph{(the
classical Bernstein-function toolbox: \sscite{feller2009introduction}, Ch.~XIII and
\sscite{schilling2012bernstein}, Ch.~1 and~3. Cited, not reproven here.
\textbf{Assignment.} The four parts are carried by four separate entries, one each.
Part~(1) is \ledger{A1} entire, in both the probability form and the general measure form
with uniqueness. Part~(2) is \ledger{A2} entire --- all three equivalences of the cited
theorem, not only the exponential one: the composition clause is what
\ref{lem:potential-kernel} inverts a Bernstein function with. Part~(3) is \ledger{A3} entire, including
the uniqueness of the triple, which Lemma~7.1 needs and which a representation theorem
without it would not supply. Part~(3) is also the \emph{bridge} between the two vocabularies
this project speaks: it is what identifies \ref{def:bernstein-function} with
\ref{def:levy-exponent}, and therefore the single edge every Lean tag in chapters 5--7 crosses,
since the development names admissibility by the representation and never by derivative signs.
A node whose conclusion is stated in $\BFz$ cannot carry a Lean tag without spending this part;
a node whose conclusion is stated in $\LE$ carries one for free.
Part~(4) is \ledger{A4} entire; it is the single external fact
Theorem~5.2 rests on, and the null-array argument that produces the pointwise limit there is
ours and is held \textbf{[T]}. None of the four carries anything about the hemigroup family:
they are statements about $\CM$ and $\BF$ as classes of functions, and every step joining
them to $\Phi_{x,y}$ is \textbf{[T]}.)}
\end{proposition}

% draft: Lemma 2.4
\begin{lemma}[vanishing lemma]\label{lem:vanishing}
\uses{def:bernstein-function,prop:bernstein-toolbox}
\lean{Hemigroup.levyExponent_eq_zero_of_eq_zero}\leanok
If $g \in \BFz$ and $g(s_0) = 0$ for some $s_0 > 0$, then $g \equiv 0$.
\statusT\quad\emph{(the Lean development states this for the \emph{representation}
$g(s) = b_0 s + \int (1 - e^{-st})\,\nu(dt)$ rather than for the derivative-sign definition
\ref{def:bernstein-function}, following the formalisation strategy; the two notions agree by
the Lévy--Khintchine theorem, which is A3 and reaches this node through
\ref{prop:bernstein-toolbox}(3). From the representation the proof is shorter than the one
below --- a sum of nonnegative terms vanishes only when each does --- and needs neither
concavity nor complete monotonicity. \texttt{\#print axioms} reduces to Lean core only.)}
\end{lemma}

\begin{proof}
\leanok
$g$ is nonnegative and concave by \ref{prop:bernstein-toolbox}(3), with
$g(\zp) = g(s_0) = 0$. A nonnegative concave function vanishing at both ends of $[0, s_0]$
vanishes on it, so $g \equiv 0$ on $(0, s_0]$ and hence $g' \equiv 0$ on $(0, s_0)$. Now $g'$
is nonincreasing, being the derivative of a concave function, and nonnegative, being
completely monotone; a nonnegative nonincreasing function that vanishes on $(0,s_0)$ vanishes
identically. So $g' \equiv 0$ on $(0,\infty)$, and with $g(\zp) = 0$ this gives $g \equiv 0$.
\end{proof}

\bigskip

The next two statements are not in the draft, which uses both facts without stating either.
They are recorded here so that the trust they carry is visible in the dependency graph rather
than buried inside proofs. Both are classical, and both are cited to the same section of
\sscite{feller2009introduction}; note that that section states the probability and the measure
forms as four separate theorems, on four different pages, so the pair a given proof needs must
be named individually.

% additive: not in the draft. Used by lem:convolution-representation, thm:main-characterization.
\begin{proposition}[continuity theorem for Laplace transforms]\label{prop:laplace-continuity}
Let $\mu_n, \mu$ be probability measures on $[0,\infty)$. Then $\mu_n \to \mu$ weakly if and
only if $\hat\mu_n(s) \to \hat\mu(s)$ for every $s > 0$. The same equivalence holds for
arbitrary measures on $[0,\infty)$, provided that in the direction from transforms to measures
the masses $\hat\mu_n(a)$ are bounded for some $a > 0$.
\statusA\quad\ledger{A5}\quad\emph{(Feller's continuity theorem;
\sscite{feller2009introduction}, Vol.~2, \S XIII.1, Theorem~2 and Theorem~2a. A cited
interface, not reproven here. \textbf{Assignment.} \ledger{A5} carries both clauses. The
boundedness hypothesis in the second is not decoration: Feller gives a counterexample on the
same page showing it cannot be dropped, and any proof here that invokes the measure form must
discharge it explicitly. The probability form carries it for free, which is why Theorem~7.3
uses that one. What \ledger{A5} does not carry is the tightness argument that supplies weak
convergence in the first place; that is \textbf{[T]}. \textbf{Not needed in Lean} (resolved
2026-08-09): the development proves the form it uses rather than citing it, as
\texttt{Hemigroup.tendsto\_integral\_of\_tendsto\_laplace}, whose \texttt{\#print axioms} reduces
to Lean core. For \emph{causal probability} measures, convergence of the transforms together with
tightness gives weak convergence; that is precisely what Theorem~7.3$(\Leftarrow)$ uses, and the
boundedness hypothesis flagged above is carried for free by the probability form, so the trap is
designed out rather than remembered. The general measure form is deliberately \emph{not}
formalised. Tightness is ours, as this clause always claimed and now checks: the Markov bound
$\mu(t>T)(1-e^{-sT}) \le 1 - \hat\mu(s)$ (\texttt{measureReal\_Ioi\_mul\_le}), uniform over the
family as $\mu_{a,b}(t>T) \le F(Bs)/(1-e^{-sT})$ (\texttt{kernel\_tail\_le}), with $F(\zp) = 0$
(\texttt{exists\_exponent\_lt}) making $s$ choosable. The node keeps \textbf{[A]} because that is
the \emph{paper's} stance, and carries no Lean tag because what is proved is one direction of the
probability case, not the statement above.)}
\end{proposition}

% additive: not in the draft. Used by lem:convolution-representation, lem:covariance-laplace,
% prop:scale-evolution, thm:sonine-conservation, lem:generator-properties.
\begin{proposition}[uniqueness of the Laplace transform]\label{prop:laplace-uniqueness}
A measure on $[0,\infty)$ is determined by its Laplace transform: if $\hat\mu(s) = \hat\rho(s)$
for all $s$ in some interval $(a,\infty)$, then $\mu = \rho$.
\statusA\quad\ledger{A6}\quad\emph{(Feller's uniqueness theorem;
\sscite{feller2009introduction}, Vol.~2, \S XIII.1, Theorem~1 and Theorem~1a. A cited
interface, not reproven here. \textbf{Assignment.} \ledger{A6} carries the statement entire,
for probability measures (Theorem~1) and for arbitrary measures (Theorem~1a). This is what
licenses every step in the article of the form ``the transforms agree, therefore the measures
agree'' --- the uniqueness clause of Lemma~4.1, and Lemma~6.1, Proposition~9.2, Theorem~9.5
and Lemma~10.3(5). It says nothing about existence, and nothing about which functions are
transforms; that is \ref{prop:bernstein-toolbox}(1). \textbf{Not needed in Lean.} The
development proves this rather than citing it: \texttt{Hemigroup.laplace\_injective}, by the
classical substitution $x = e^{-t}$, which carries a causal measure onto $[0,1]$, turns
$\hat\mu(n)$ into the $n$th moment of the image, and closes with the Weierstrass approximation
theorem. \texttt{\#print axioms} reduces to Lean core, so A6 never enters
\texttt{trust-boundary.txt} and nothing proved in Lean rests on it. The node keeps \textbf{[A]}
because that is the \emph{paper's} stance, and it carries no Lean tag because what is proved is
weaker than what is stated: agreement on all of $[0,\infty)$, not on a tail $(a,\infty)$. That
gap is not a formality --- equality of the moments from some $n_0$ onwards does not by itself
pin down an atom at the origin of the transported measure, and the argument that ours has none
is not formalised.)}
\end{proposition}

\bigskip

The remaining four statements are additive too, and they are the \emph{formalisation's}
vocabulary rather than the draft's. \ref{def:levy-exponent} names the object the Lean
development is written in; the three after it are the restricted cases of \ref{prop:laplace-continuity}
and \ref{prop:laplace-uniqueness} that it proves outright, recorded as \textbf{[T]} nodes beneath
the \textbf{[A]} ones they refine so that the graph shows what actually rests on \ledger{A5} and
\ledger{A6}. Nothing formalised does.

% additive: not in the draft. The representation as a named class -- what the Lean
% development speaks. Used by lem:vanishing, thm:increments-bernstein, prop:canonical-gauge,
% lem:selfdecomposable-exponents, thm:main-construction.
\begin{definition}[Lévy exponent]\label{def:levy-exponent}
\uses{prop:bernstein-toolbox}
\lean{Hemigroup.levyExponent}\leanok
A function $g : [0,\infty) \to [0,\infty]$ is a \emph{Lévy exponent} if
\[
  g(s) = b_0\,s + \int_0^\infty \bigl(1 - e^{-st}\bigr)\,\nu(dt),
  \qquad b_0 \ge 0,
\]
for a drift $b_0$ and a positive Borel measure $\nu$ on $(0,\infty)$ --- with no killing term,
so that $g(0) = 0$. We write $g \in \LE$, and call $(b_0, \nu)$ its \emph{triple}.
\statusT\quad\emph{(A \emph{name} for \ref{prop:bernstein-toolbox}(3)'s right-hand side, not a
new fact: by that part, $\LE = \BFz$ as classes of functions, with the triple unique. Both names
are kept because the two documents need different ones. The paper argues in $\BFz$, where the
draft's proofs live; the Lean development argues in $\LE$ and never defines
\ref{def:completely-monotone} at all --- the reason is \texttt{blueprint/DESIGN-formalization-strategy.md},
and the payoff is that the Lévy--Khintchine theorem (\ledger{A3}) is crossed once, here, instead
of inside every proof. Stated in Lean as \texttt{Hemigroup.levyExponent}, valued in
$[0,\infty]$ so that no integrability hypothesis is carried; the integrability condition
$\int (1 \wedge t)\,\nu(dt) < \infty$ of \ref{prop:bernstein-toolbox}(3) is equivalent, for
causal $\nu$, to finiteness of $g$ at a single $s > 0$, and appears in the development as that
hypothesis instead. \ref{lem:vanishing} and \ref{lem:selfdecomposable-increment}
are both proved directly from this definition
(\texttt{levyExponent\_eq\_zero\_of\_eq\_zero}, \texttt{levyExponentD\_increment}).)}
\end{definition}

% additive: not in the draft. The [T] refinement of prop:laplace-uniqueness.
\begin{proposition}[uniqueness of the Laplace transform, causal case]\label{prop:laplace-uniqueness-causal}
\uses{prop:laplace-uniqueness}
\lean{Hemigroup.laplace_injective}\leanok
Let $\mu, \rho$ be finite measures on $\RR$ carried by $[0,\infty)$. If
$\hat\mu(s) = \hat\rho(s)$ for every $s \ge 0$, then $\mu = \rho$.
\statusT\quad\emph{(the case of \ref{prop:laplace-uniqueness} that this article's own arguments
use, and it is \emph{proved}: \texttt{Hemigroup.laplace\_injective},
\texttt{\#print axioms} reducing to Lean core. It is weaker than the cited statement in exactly
one respect --- agreement is assumed on all of $[0,\infty)$, not on a tail $(a,\infty)$ --- and
that gap is not a formality, which is why \ref{prop:laplace-uniqueness} keeps its \ledger{A6}
tag and this node does not replace it. Every step in the formalised part of the article of the
form ``the transforms agree, therefore the measures agree'' has both measures carried by
$[0,\infty)$ and both transforms known on all of it, so it cites this node and not that one; the
consequence is that \ledger{A6} does not appear in \texttt{blueprint/trust-boundary.txt}.)}

\emph{\textbf{Extended 2026-08-11 to measures that are not finite}, which is what Chapter~9
needs: $\kappa^{(x)}$, $\ell^{(x)}$ and Lebesgue measure are none of them finite, and
\ref{thm:sonine-conservation} compares them. That is \ref{prop:laplace-uniqueness-sigma-finite}
below. The finite case remains the one almost every caller uses, so it keeps this node and its
tag.}
\end{proposition}

\begin{proof}
\leanok
Substitute $x = e^{-t}$: the map is a measurable embedding carrying a measure on $[0,\infty)$
to one on $(0,1]$, and it turns $\hat\mu(n)$ into the $n$th moment of the image, for every
$n \in \NN$. Two finite measures on the compact $[0,1]$ with the same moments agree, by the
Weierstrass approximation theorem. Transporting back along the embedding loses nothing, so
$\mu = \rho$.
\end{proof}

% additive: not in the draft. The unbounded-measure refinement of A6, split from
% prop:laplace-uniqueness-causal because its hypothesis is genuinely different -- not finiteness
% of the measure but convergence of the transform. Numbered 2.13.
\begin{proposition}[uniqueness without a finiteness assumption]\label{prop:laplace-uniqueness-sigma-finite}
\uses{prop:laplace-uniqueness-causal}
\lean{Hemigroup.laplaceL_injective_of_ne_top}\leanok
Let $\mu, \rho$ be measures on $\RR$ carried by $[0,\infty)$, not assumed finite, and let
$s_0 \in \RR$ be such that $\hat\mu(s_0) < \infty$. If $\hat\mu(s) = \hat\rho(s)$ for every
$s \ge s_0$, then $\mu = \rho$.
\statusT\quad\emph{(\textbf{Proved}, as \texttt{Hemigroup.laplaceL\_injective\_of\_ne\_top},
reducing to Lean core --- so \ledger{A6} stays off the trust boundary in its general form as
well as its restricted one. The hypothesis is worth reading carefully, because it is weaker than
the one the chapter that needs this was expected to supply: what is assumed is not that the
measures are locally finite but that \emph{the transform converges at a single point}. Local
finiteness is then a consequence, \ref{lem:laplace-local-finiteness}, not an assumption.}

\emph{The proof is the classical damping trick and nothing more. $e^{-s_0 t}\mu(dt)$ is finite
precisely because the transform converges at $s_0$; its transform at $s$ is $\hat\mu(s_0 + s)$,
so \ref{prop:laplace-uniqueness-causal} applies to the damped pair; and the damping is undone by
multiplying the density back, which is legitimate because $e^{-s_0t}$ is everywhere positive and
finite. Nothing here needs the tail-agreement form of \ledger{A6} that
\ref{prop:laplace-uniqueness-causal} declines to claim.)}
\end{proposition}

\begin{proof}
\leanok
The measure $\mu_0(dt) := e^{-s_0 t}\,\mu(dt)$ is finite, its mass being $\hat\mu(s_0) < \infty$,
and is carried by $[0,\infty)$, being absolutely continuous with respect to $\mu$; likewise
$\rho_0$, whose mass is $\hat\rho(s_0) = \hat\mu(s_0)$. For $s \ge 0$,
$\hat{\mu_0}(s) = \hat\mu(s_0 + s) = \hat\rho(s_0 + s) = \hat{\rho_0}(s)$, so $\mu_0 = \rho_0$ by
\ref{prop:laplace-uniqueness-causal}. Since $e^{-s_0t}\,e^{s_0t} = 1$ pointwise, multiplying both
by the density $e^{s_0 t}$ recovers $\mu = \rho$.
\end{proof}

\begin{lemma}[convergence forces local finiteness]\label{lem:laplace-local-finiteness}
\uses{prop:laplace-uniqueness-sigma-finite}
\lean{Hemigroup.measure_Icc_ne_top_of_laplaceL_ne_top}\leanok
Let $\mu$ be carried by $[0,\infty)$ with $\hat\mu(s_0) < \infty$ for some $s_0 > 0$. Then
$\mu([0,T]) < \infty$ for every $T$.
\statusT\quad\emph{(\textbf{Proved}, as
\texttt{Hemigroup.measure\_Icc\_ne\_top\_of\_laplaceL\_ne\_top}. Stated because it is what makes
the hypothesis of \ref{prop:laplace-uniqueness-sigma-finite} the natural one: ``locally finite''
would have been an extra assumption, and it is instead a consequence.)}
\end{lemma}

\begin{proof}
\leanok
On $[0,T]$ the integrand satisfies $1 \le e^{s_0 T} e^{-s_0 t}$, so
$\mu([0,T]) \le e^{s_0T}\hat\mu(s_0) < \infty$.
\end{proof}

% additive: not in the draft. The [T] refinement of prop:laplace-continuity.
\begin{proposition}[continuity theorem, causal probability case]\label{prop:laplace-continuity-causal}
\uses{prop:laplace-continuity,prop:laplace-uniqueness-causal}
\lean{Hemigroup.tendsto_integral_of_tendsto_laplace}\leanok
Let $\mu_n, \mu$ be probability measures on $\RR$ carried by $[0,\infty)$, and suppose
\begin{enumerate}
\item $\hat\mu_n(s) \to \hat\mu(s)$ for every $s \ge 0$; and
\item the family is \emph{tight}: for every $\eta > 0$ there is a $T$ with
$\mu_n\bigl((T,\infty)\bigr) \le \eta$ for all $n$, and $\mu\bigl((T,\infty)\bigr) \le \eta$.
\end{enumerate}
Then $\mu_n \to \mu$ weakly.
\statusT\quad\emph{(the direction and the case of \ref{prop:laplace-continuity} that
\ref{thm:main-construction} uses --- transforms to measures, probability form --- and it is
\emph{proved}: \texttt{Hemigroup.tendsto\_integral\_of\_tendsto\_laplace},
\texttt{\#print axioms} reducing to Lean core, so \ledger{A5} does not appear in
\texttt{blueprint/trust-boundary.txt} either. Tightness is a \emph{hypothesis} here rather than a
conclusion, which is the honest shape: \ref{prop:laplace-continuity}'s assignment clause always
held the tightness argument \textbf{[T]}, and \ref{lem:transform-tightness} is where it is
discharged. Feller's boundedness hypothesis, the classical trap of the general form, is carried
for free by the probability form and so cannot be forgotten. The general measure form is
deliberately not formalised.)}
\end{proposition}

\begin{proof}
\leanok
Push forward along $x = e^{-t}$ onto the compact $[0,1]$, as in
\ref{prop:laplace-uniqueness-causal}. There, convergence of the Laplace transforms at the
naturals \emph{is} convergence of the moments, so the images converge weakly by Weierstrass; no
compactness argument is needed, because the carrier is already compact.

Transporting back is where tightness is spent. A bounded continuous $f$ on $[0,\infty)$ pulls
back to $f \circ (-\log)$, which is unbounded at the origin, so it is replaced by
$f \circ \bigl(\min(-\log, T)\bigr)$, bounded and continuous on $[0,1]$ and agreeing with it off
$(T,\infty)$; the two integrals differ by at most $2\|f\|_\infty\,\mu_n\bigl((T,\infty)\bigr)$,
which hypothesis (2) makes uniformly small. An $\varepsilon/3$ argument concludes.
\end{proof}

% additive: not in the draft. The tightness half prop:laplace-continuity's assignment clause
% always held [T] but never stated. Used by thm:main-construction.
\begin{lemma}[tail bound from the transform]\label{lem:transform-tightness}
\lean{Hemigroup.measureReal_Ioi_mul_le}\leanok
Let $\mu$ be a probability measure on $\RR$ carried by $[0,\infty)$. Then for all $s \ge 0$ and
all $T$,
\[
  \mu\bigl((T,\infty)\bigr)\,\bigl(1 - e^{-sT}\bigr) \;\le\; 1 - \hat\mu(s).
\]
Consequently a family of such measures whose transforms satisfy
$\sup_n \bigl(1 - \hat\mu_n(s)\bigr) \to 0$ as $s \to \zp$ is tight.
\statusT\quad\emph{(this is the tightness argument \ref{prop:laplace-continuity} always claimed
as ours, now stated rather than left inside a proof. \texttt{Hemigroup.measureReal\_Ioi\_mul\_le},
\texttt{\#print axioms} reducing to Lean core; the divided form is
\texttt{measureReal\_Ioi\_le\_div}. No hypothesis on $T$ is needed --- causality supplies the
sign that $T \ge 0$ would have. Its use in \ref{thm:main-construction} is
\texttt{kernel\_tail\_le}, which makes the bound uniform over the family as
$\mu_{a,b}\bigl((T,\infty)\bigr) \le F(Bs)/(1 - e^{-sT})$ for $b \le B$, and what makes that
bite is $F(\zp) = 0$; \texttt{isTightMeasureSet\_kernel} repackages it in Mathlib's own
vocabulary, which nothing here now consumes but any other consumer would.)}
\end{lemma}

\begin{proof}
\leanok
Pointwise: $1 - \hat\mu(s) = \int \bigl(1 - e^{-st}\bigr)\,\mu(dt)$, the integrand is at least
$1 - e^{-sT}$ on $(T,\infty)$, and it is nonnegative elsewhere because $\mu$ is carried by
$[0,\infty)$ and $s \ge 0$. Dropping the complement gives the bound.

For the consequence, fix $\eta > 0$. Choose $s > 0$ with
$\sup_n \bigl(1 - \hat\mu_n(s)\bigr) \le \eta/2$, then $T$ large enough that
$1 - e^{-sT} \ge 1/2$; the displayed bound gives $\mu_n\bigl((T,\infty)\bigr) \le \eta$ for
every $n$.
\end{proof}


\chapter{Convolution representation}
% !TEX root = ../print.tex
% Draft \S4 --- Convolution representation
%
% Source: draft/hemigroup-causal-scale-space-kernels.md, lines 93-119.
% Chapter 4 of the blueprint; the shared counter puts every number here on the draft's own.
%
% The first two \textbf{[T]} nodes carrying real proofs. Both are proved from scratch: this is
% the Wendel step (translation-invariant operators on $L^1$ are convolutions), and \S1 and \S11
% of the draft name it as such, but the argument here is self-contained -- an approximate
% identity plus a weak-* subnet -- and cites nothing. \sscite{wendel1952left} is therefore
% mentioned in the remark for provenance, NOT as a \ledger{} interface: nothing here is taken
% on trust.
%
% Banach--Alaoglu and Riesz--Markov (the weak-* extraction) are likewise held [T]. They are
% deliberately not ledger entries -- Mathlib has both, so the Lean development can prove this
% rather than assume it, which is exactly the verified-core/axiomatized-analysis split working
% as intended.

% draft: Lemma 4.1
\begin{lemma}[representation]\label{lem:convolution-representation}
\uses{def:cascade-family,prop:laplace-uniqueness}
Under (A1)--(A5), for each pair $x \le y$ there is a unique probability measure $\mu_{x,y}$ on
$[0,\infty)$ with
\[
  \Phi_{x,y} f = \mu_{x,y} * f, \qquad f \in X,
\]
and conversely every such measure defines an operator satisfying (A1)--(A5).
\statusT
\end{lemma}

\begin{proof}
Fix $(x,y)$ and write $\Phi = \Phi_{x,y}$. Choose an approximate identity
$\rho_\varepsilon \in X_+$ with $\int \rho_\varepsilon = 1$ and
$\supp \rho_\varepsilon \subseteq [0,\varepsilon]$. By (A4) and (A5),
$h_\varepsilon := \Phi \rho_\varepsilon \in X_+$ with $\int h_\varepsilon = 1$; by (A3),
$h_\varepsilon$ vanishes a.e.\ on $(-\infty, 0)$. Regarding $h_\varepsilon\,dt$ as elements of
the unit ball of $M(\RR) = C_0(\RR)^*$, extract a subnet converging weak-$*$ to some
$\mu \ge 0$ with $\|\mu\| \le 1$ and $\supp \mu \subseteq [0,\infty)$.

For $f, g \in X$ we have $f * g = \int \shift_r g\, f(r)\,dr$ as a Bochner integral, so
(A1)--(A2) give $\Phi(f * g) = f * \Phi g$. Hence for $f \in C_c(\RR)$,
\[
  \Phi(f * \rho_\varepsilon) = f * h_\varepsilon
  \xrightarrow{\ \varepsilon \to 0\ } f * \mu \quad \text{pointwise},
\]
using $f(t - \cdot) \in C_0(\RR)$ and weak-$*$ convergence, while
$\Phi(f * \rho_\varepsilon) \to \Phi f$ in $X$ by (A1). Therefore $\Phi f = \mu * f$ for
$f \in C_c$, and by density for all $f \in X$.

Finally, for $f \in X_+$, Tonelli gives $\int \mu * f = \|\mu\|\int f$, so (A5) forces
$\|\mu\| = 1$. Uniqueness holds because a finite measure on $[0,\infty)$ is determined by its
Laplace transform (\ref{prop:laplace-uniqueness}) and $\widehat{\mu * f} = \hat\mu\,\hat f$
with $\hat f \not\equiv 0$ for some $f \in X_+$. For the converse, a probability measure
$\mu$ on $[0,\infty)$ makes $f \mapsto \mu * f$ bounded with norm $\le 1$ by Young's
inequality (A1), translation-commuting because convolution is (A2), support-preserving on
$[t_0,\infty)$ because $\supp\mu \subseteq [0,\infty)$ (A3), positive (A4), and mass-preserving
by Tonelli (A5).
\end{proof}

% draft: Lemma 4.2
\begin{lemma}[continuity of transforms]\label{lem:transform-continuity}
\uses{def:cascade-family,lem:convolution-representation}
Under (A1)--(A5) and (A7), the function
\[
  g_{x,y}(s) := -\log \hat\mu_{x,y}(s) \in [0,\infty), \qquad s \ge 0,
\]
is well defined, and $(x,y) \mapsto g_{x,y}(s)$ is continuous for each $s \ge 0$; moreover
$g_{x,y}(\zp) = 0$ and $s \mapsto g_{x,y}(s)$ is continuous.
\statusT
\end{lemma}

\begin{proof}
$\hat\mu(s) \in (0, 1]$ for a probability measure on $[0,\infty)$, so $g$ is well defined and
nonnegative; $\hat\mu(s) \to 1$ as $s \downarrow 0$ by dominated convergence, and
$s \mapsto \hat\mu(s)$ is continuous.

For continuity in $(x,y)$: fix a probability density $\rho \in X_+$ supported in $[0,\infty)$,
so that $\hat\rho(s) > 0$. If $(x_n, y_n) \to (x,y)$, then by (A7)
$\mu_{x_n,y_n} * \rho \to \mu_{x,y} * \rho$ in $X$; since all functions involved are supported
in $[0,\infty)$ and $|e^{-st}| \le 1$ there,
\[
  \bigl|\hat\mu_{x_n,y_n}(s)\,\hat\rho(s) - \hat\mu_{x,y}(s)\,\hat\rho(s)\bigr|
  \le \|\mu_{x_n,y_n} * \rho - \mu_{x,y} * \rho\|_1 \longrightarrow 0,
\]
and dividing by $\hat\rho(s) > 0$ gives the claim.
\end{proof}

% Not in the draft as a numbered item; the observation is made in its \S1 and \S11 prose.
\begin{remark}[what the representation is, and what it is not]\label{rem:wendel}
\uses{lem:convolution-representation}
\ref{lem:convolution-representation} is the Wendel step
(\sscite{wendel1952left}: the translation-invariant operators on $L^1$ of a group are exactly
convolution by a measure), specialised to $\RR$ and sharpened by (A3)--(A5) to a probability
measure carried by $[0,\infty)$. It is proved here rather than cited, so nothing in this
chapter rests on the trust boundary: the weak-$*$ extraction is Banach--Alaoglu together with
the Riesz--Markov identification $M(\RR) = C_0(\RR)^*$, both of which the formalisation can
discharge outright.

The multiplicative-group counterpart of this step, for the scale variable rather than time,
is what Chapter~11 has to build by hand; there is no Wendel theorem to appeal to once the
group is $(0,\infty)$ under multiplication and the operator is required only to be covariant.
\end{remark}


\chapter{The cascade: additivity and infinite divisibility}
% !TEX root = ../print.tex
% Draft \S5 --- The cascade
%
% Source: draft/hemigroup-causal-scale-space-kernels.md, lines 121-160.
% This file is chapter 5 --- The cascade of the blueprint; the shared counter makes every number here
% agree with the draft's own (see content.tex).
%
% Planned nodes: lem:additivity [T], thm:increments-bernstein [T], cor:strict-monotonicity [T]
%
% No nodes yet. Each statement carries \label{kind:name}, \lean{...} with
% \leanok/\notready, \uses{...}, \notes{slug}, a \statusT or \statusA line, and a
% "% draft: <kind> N.M" provenance comment; an \statusA node must also declare what its
% citation carries (\textbf{Assignment.}).


\chapter{Scale covariance: the similarity form}
% !TEX root = ../print.tex
% Draft \S6 --- Scale covariance: the similarity form
%
% Source: draft/hemigroup-causal-scale-space-kernels.md, lines 161-196.
% Chapter 6 of the blueprint; the shared counter puts every number here on the draft's own.
%
% The chapter that turns "there is some action of the scaling group" into a coordinate. Its
% two rigidity results are proved in the draft as single dense paragraphs -- Lemma 6.2 argues
% four numbered claims in one, with the key monotonicity step inside an em-dash aside -- and
% they are broken out here, statement by statement, with no change of content. That is the
% one place in this chapter where the blueprint should read differently from the draft: this
% is the argument a formalisation has to follow line by line, and the draft's compression is
% the enemy of exactly that. The formalisation note singles out this gauge construction as
% the most unglamorous stretch of the Lean ladder (its M5); the proofs below are written to
% be followed rather than admired.
%
% Equations (6.1) and (6.2) keep the draft's hand-written tags, since its later prose refers
% to them by number.
%
% VOCABULARY (2026-08-10), the companion to the note at the head of 05-cascade.tex. Only one
% conclusion in this chapter names a class of functions -- prop:canonical-gauge's F in BF_0 --
% and it is now also stated in LE (def:levy-exponent), which is what a Lean tag can point at.
%
% The proofs are a second matter and a smaller one. Two steps here reach for the derivative-sign
% picture where they do not have to: lem:action-rigidity(3) argues that a nonzero member of BF_0
% is strictly increasing "because its derivative is completely monotone", and the same fact is
% used again in prop:canonical-gauge. Both are cor:strict-monotonicity, whose Lean route runs
% lem:vanishing on the representation and never differentiates -- see that node. The prose is
% left as the draft has it, since the paper argues in BF_0; the formalisation notes say which
% substitute to use, so that a Lean pass does not rediscover the question mid-proof.

% draft: Lemma 6.1
\begin{lemma}[Laplace form of covariance]\label{lem:covariance-laplace}
\uses{def:cascade-family,lem:convolution-representation,lem:additivity,prop:laplace-uniqueness}
\lean{Hemigroup.CascadeCore.covariance_laplace}\leanok
Under (A1)--(A5), axiom (A8) is equivalent to
\[
  \dil_\sigma\,\mu_{x,y} = \mu_{\Sact_\sigma x,\, \Sact_\sigma y}
  \qquad \Longleftrightarrow \qquad
  g_{\Sact_\sigma x,\, \Sact_\sigma y}(s) = g_{x,y}(\sigma s),
\]
and in particular
\[
  G(\Sact_\sigma x,\, s) = G(x,\, \sigma s)
  \qquad \text{for all } x \ge 0,\ \sigma, s > 0. \tag{6.1}
\]
\statusT\quad\emph{the ``is equivalent to'' is carried by two declarations: (A8)
$\Rightarrow$ identities is \texttt{covariance\_laplace}, and identities $\Rightarrow$ (A8) is
\texttt{isScaleCovariant\_of\_repr\_map} (\texttt{Covariance.lean}, fidelity review R18).}
\end{lemma}

\begin{proof}
\leanok
The compatibility $\dil_\sigma(\mu * f) = (\dil_\sigma\mu) * (\dil_\sigma f)$ turns (A8) into
the statement that $\dil_\sigma\mu_{x,y}$ and $\mu_{\Sact_\sigma x, \Sact_\sigma y}$ represent
the same operator, and the uniqueness clause of \ref{lem:convolution-representation} makes
that an identity of measures; this is the first equivalence. For the second,
$\widehat{\dil_\sigma\mu}(s) = \hat\mu(\sigma s)$, so taking $-\log$ gives
$g_{\Sact_\sigma x, \Sact_\sigma y}(s) = g_{x,y}(\sigma s)$. Equation (6.1) is the case
$x = 0$ of that identity together with $\Sact_\sigma 0 = 0$, which holds because an increasing
bijection of $[0,\infty)$ fixes $0$; $G$ is $g_{0,\cdot}$ by \ref{lem:additivity}.
\end{proof}

% draft: Lemma 6.2
\begin{lemma}[rigidity of the action]\label{lem:action-rigidity}
\uses{def:cascade-family,lem:covariance-laplace,cor:strict-monotonicity,lem:additivity}
\lean{Hemigroup.CascadeCore.action_rigidity}\leanok
Assume (A1)--(A8) and (ND). Then:
\begin{enumerate}
\item For each $\sigma$, $\Sact_\sigma$ is uniquely determined by (6.1).
\item $\Sact_\sigma \Sact_\tau = \Sact_{\sigma\tau}$ for all $\sigma, \tau > 0$, and
$\Sact_1 = \Id$.
\item $\sigma \mapsto \Sact_\sigma x$ is continuous and strictly increasing for each $x > 0$.
\item The action has no fixed point in $(0,\infty)$: if $\Sact_\sigma x^* = x^*$ for all
$\sigma$, then $x^* = 0$.
\end{enumerate}
\statusT\quad\emph{(three of the four clauses are injectivity of $G(\cdot,s)$ --- that is,
\ref{cor:strict-monotonicity} --- and nothing else. (1) \emph{is} that injectivity
(\texttt{eq\_of\_G\_eq}); (2) is two applications of (6.1) on top of it (\texttt{S\_comp},
\texttt{S\_one}); (4) adds only that $G(x,\cdot)$ is continuous at the origin with value $0$,
which turns ``constant on $(0,\infty)$'' into ``identically zero'' and so contradicts
\ref{cor:strict-monotonicity} above $x = 0$ (\texttt{eq\_zero\_of\_fixed}).}

\emph{What (3) needed, and the development did not have, was strict monotonicity in the
\emph{other} variable: the proof below reads it off the sign of a Bernstein function's
derivative. The substitute is \texttt{laplace\_strictAnti}, one line of measure theory --- the
support of $e^{-st} - e^{-s't}$ is exactly $\{t \ne 0\}$, and that set carries positive mass
precisely because the kernel is not $\delta_0$, which is (ND). No derivative and no Bernstein
class appears. With it, (3)'s continuity half does not need the inverse function the proof below
invokes: \texttt{StrictMonoOn.continuousAt\_of\_exists\_between} asks only for orbit points
squeezed between a given bound and $\Sact_{\sigma_0} x$, and continuity of
$\sigma \mapsto G(x,\sigma)$ supplies them. That also avoids having to show the range of
$G(\cdot,s_0)$ is a neighbourhood, which is the step the phrase ``on its range'' hides.)}
\end{lemma}

\begin{proof}
\leanok
\emph{(1) Uniqueness.} If $\Sact$ and $\Sact'$ both satisfy (6.1) then
$G(\Sact_\sigma x, s) = G(x,\sigma s) = G(\Sact'_\sigma x, s)$ for every $s > 0$. By
\ref{cor:strict-monotonicity}, $G(\cdot, s)$ is strictly increasing, hence injective, so
$\Sact_\sigma x = \Sact'_\sigma x$.

\emph{(2) The group law.} Applying (6.1) twice,
\[
  G(\Sact_\sigma \Sact_\tau x, s) = G(\Sact_\tau x, \sigma s) = G(x, \sigma\tau s)
  = G(\Sact_{\sigma\tau} x, s),
\]
and injectivity of $G(\cdot,s)$ gives $\Sact_\sigma \Sact_\tau = \Sact_{\sigma\tau}$. Taking
$\sigma = 1$ in (6.1) gives $G(\Sact_1 x, s) = G(x,s)$, whence $\Sact_1 = \Id$.

\emph{(3) Continuity and strict monotonicity in $\sigma$.} Fix $x > 0$ and any $s_0 > 0$. By
(6.1), $G(\Sact_\sigma x, s_0) = G(x, \sigma s_0) = g_{0,x}(\sigma s_0)$.

The map $\sigma \mapsto g_{0,x}(\sigma s_0)$ is continuous, and it is \emph{strictly}
increasing: $g_{0,x}$ is a nonzero member of $\BFz$ by \ref{cor:strict-monotonicity}, and a
nonzero member of $\BFz$ is strictly increasing, because its derivative is completely monotone
and not identically zero, hence strictly positive on $(0,\infty)$ --- a completely monotone
function vanishing at one point vanishes beyond it, and vanishing on an interval forces the
representing measure, and so the function, to be zero.

The map $G(\cdot, s_0)$ is continuous (\ref{lem:additivity}) and strictly increasing
(\ref{cor:strict-monotonicity}), so it has a continuous strictly increasing inverse on its
range. Composing,
$\Sact_\sigma x = G(\cdot,s_0)^{-1}\bigl(g_{0,x}(\sigma s_0)\bigr)$ is continuous and strictly
increasing in $\sigma$.

\emph{(4) No interior fixed point.} Suppose $\Sact_\sigma x^* = x^*$ for all $\sigma > 0$.
Then (6.1) gives $G(x^*, s) = G(x^*, \sigma s)$ for all $\sigma, s > 0$, so $G(x^*, \cdot)$ is
constant on $(0,\infty)$; its value is $G(x^*, \zp) = 0$ by \ref{lem:additivity}. Hence
$g_{0,x^*} \equiv 0$, i.e.\ $\mu_{0,x^*} = \delta_0$ and $\Phi_{0,x^*} = \Id$. By (ND) this
forces $x^* = 0$.
\end{proof}

% draft: Proposition 6.3
\begin{proposition}[orbit coordinate; the gauge]\label{prop:canonical-gauge}
\uses{lem:action-rigidity,lem:covariance-laplace,cor:strict-monotonicity,thm:increments-bernstein,def:bernstein-function,def:levy-exponent}
\lean{Hemigroup.CascadeCore.canonical_gauge}\leanok
Assume (A1)--(A8) and (ND). Then $c(\sigma) := \Sact_\sigma 1$ is a continuous strictly
increasing bijection of $(0,\infty)$ onto $(0,\infty)$ with $c(1) = 1$. Consequently
\[
  \chi := c^{-1} : (0,\infty) \to (0,\infty), \qquad \chi(0) := 0,
\]
is an increasing bijection --- \emph{the canonical gauge} --- satisfying
$\chi(\Sact_\sigma x) = \sigma\,\chi(x)$, and
\[
  G(x, s) = F\bigl(\chi(x)\, s\bigr), \qquad
  F := G(1, \cdot) \in \BFz,\ F \not\equiv 0. \tag{6.2}
\]
Equivalently $F \in \LE$, with a drift $b_0 \ge 0$ and a Lévy measure $\nu$ on $(0,\infty)$.
\statusT\quad\emph{(the $\LE$ reading is the same statement, by \ref{prop:bernstein-toolbox}(3),
and is the one \ref{thm:main-analysis} needs: it is what feeds
\ref{lem:selfdecomposable-exponents} and, through it, the form (7.1). $F$'s membership comes
straight from \ref{thm:increments-bernstein}, so nothing new is asserted here --- the class is
carried along, not re-established. \textbf{Formalised}, and the interface question the
formalisation note flagged (M5 of the ladder) is settled the way it warned it had to be:
\texttt{canonical\_gauge} asserts $\chi$ together with the conjugation
$\chi(\Sact_\sigma x) = \sigma\,\chi(x)$ and the similarity form $G(x,s) = G(1,\chi(x)\,s)$,
not $\chi$ alone. That bundle is what \ref{prop:main-uniqueness} declined to guess at.}

\emph{Only surjectivity of $c$ had to be proved; continuity and strict monotonicity are
\ref{lem:action-rigidity}(3) and $c(1) = 1$ is the group law, so the injectivity argument below
--- the iteration along $\rho^{\pm n}$ --- is not needed in Lean at all. Surjectivity is done
without limits. The orbit $R = c((0,\infty))$ is invariant under every $\Sact_\rho$, which
\emph{is} the group law; were it bounded above, its supremum would be fixed by the action ---
nothing on the orbit can be carried above it, and nothing below it either, since
$\Sact_{\rho^{-1}}$ would then carry it down --- and \ref{lem:action-rigidity}(4) forbids a fixed
point above the origin, while $1 \in R$. Bounded below is the same argument read downwards, where
a positive infimum is itself the contradiction. The intermediate value theorem then finishes.)}
\end{proposition}

\begin{proof}
\leanok
\emph{$c$ is continuous} by \ref{lem:action-rigidity}(3), and $c(1) = \Sact_1 1 = 1$ by
\ref{lem:action-rigidity}(2).

\emph{$c$ is injective.} Suppose $c(\sigma) = c(\tau)$ with $\rho := \sigma/\tau \ne 1$. The
group law \ref{lem:action-rigidity}(2) gives $\Sact_\rho 1 = 1$, so by (6.1)
\[
  F(s) = G(1,s) = G(\Sact_\rho 1, s) = G(1, \rho s) = F(\rho s)
  \qquad \text{for all } s > 0 .
\]
Iterating along $\rho^{\pm n}$ gives $F(s) = F(\rho^n s)$ for every $n \in \ZZ$. Replacing
$\rho$ by $\rho^{-1}$ if necessary we may take $\rho < 1$, and then $\rho^n s \to 0$, so
$F(s) = F(\zp) = 0$ for every $s$ by \ref{lem:additivity}. Thus $F \equiv 0$, which
contradicts \ref{cor:strict-monotonicity} under (ND).

\emph{$c$ is increasing.} Being continuous and injective on the connected set $(0,\infty)$,
$c$ is strictly monotone, so it suffices to exclude the decreasing case. Suppose $c$ were
decreasing and fix $\sigma > 1$, so that $\Sact_\sigma 1 < 1$. Since $G(\cdot, s)$ is strictly
increasing (\ref{cor:strict-monotonicity}),
\[
  G(\Sact_\sigma 1, s) < G(1,s) = F(s) \qquad \text{for every } s > 0 .
\]
But (6.1) evaluates the left-hand side as $G(1,\sigma s) = F(\sigma s)$, and $F$ is
nondecreasing, being a Bernstein function, so $F(\sigma s) \ge F(s)$. The two are
incompatible, so $c$ is increasing.

\emph{$c$ is onto $(0,\infty)$.} Its limits as $\sigma \to \zp$ and $\sigma \to \infty$ exist
in $[0,\infty]$ by monotonicity, and each is a fixed point of the action or a boundary point;
by \ref{lem:action-rigidity}(4) there is no fixed point in $(0,\infty)$, so the limits are $0$
and $\infty$. With continuity, $c$ is a bijection onto $(0,\infty)$.

\emph{The gauge relation.} By the group law, $\Sact_\sigma x = \Sact_\sigma \Sact_{\chi(x)} 1
= \Sact_{\sigma\chi(x)} 1$, so $\chi(\Sact_\sigma x) = \sigma\chi(x)$.

\emph{The similarity form.} For $x > 0$, writing $x = \Sact_{\chi(x)} 1$ and using (6.1),
\[
  G(x, s) = G\bigl(\Sact_{\chi(x)} 1,\, s\bigr) = G\bigl(1,\, \chi(x)\, s\bigr)
  = F\bigl(\chi(x)\,s\bigr),
\]
and $F = g_{0,1} \in \BFz$ by \ref{thm:increments-bernstein}, with $F \not\equiv 0$ by (ND)
through \ref{cor:strict-monotonicity}. The case $x = 0$ is $G(0,\cdot) = 0 = F(0)$.
\end{proof}

% draft: Remark 6.4
\begin{remark}[gauge freedom]\label{rem:gauge-freedom}
\uses{prop:canonical-gauge}
In the hemigroup setting the scale coordinate carries no intrinsic parametrization: replacing
$x$ by any increasing bijection of $[0,\infty)$ yields an equivalent family.
\ref{prop:canonical-gauge} fixes the \emph{canonical gauge} $\tilde x := \chi(x)$, in which
$\Sact_\sigma$ acts by multiplication, scale carries the dimension of time, and
$G(x,s) = F(\tilde x\, s)$.

The powers $\tilde x \mapsto \tilde x^{\,\alpha}$ parametrize the other homogeneous gauges;
the \emph{parabolic gauge} $\alpha = 1/2$ is the natural one for the diffusive memory-line
realization and is used in the sequel.
\end{remark}


\chapter{Self-decomposability and the main theorem}
% !TEX root = ../print.tex
% Draft \S7 --- Self-decomposability and the main theorem
%
% Source: draft/hemigroup-causal-scale-space-kernels.md, lines 198-253.
% This file is chapter 7 --- Self-decomposability and the main theorem of the blueprint; the shared counter makes every number here
% agree with the draft's own (see content.tex).
%
% Planned nodes: lem:selfdecomposable-exponents [T], thm:main-characterization [T], cor:semigroup-case [T]; Remarks 7.2, 7.5
%
% No nodes yet. Each statement carries \label{kind:name}, \lean{...} with
% \leanok/\notready, \uses{...}, \notes{slug}, a \statusT or \statusA line, and a
% "% draft: <kind> N.M" provenance comment; an \statusA node must also declare what its
% citation carries (\textbf{Assignment.}).


\chapter{Examples and moments}
% !TEX root = ../print.tex
% Draft \S8 --- Examples and moments
%
% Source: draft/hemigroup-causal-scale-space-kernels.md, lines 254-278.
% Chapter 8 of the blueprint; the shared counter puts every number here on the draft's own.
%
% The draft calls 8.1-8.3 Examples. They are transcribed as PROPOSITIONS, because they are
% load-bearing rather than illustrative: \ref{thm:locality} classifies the Bessel-K family as
% the order-2 local case, \ref{prop:local-ladder} builds on the Gamma family, and \S\S13-14
% (outside this transcription) rest on both. A \uses{} edge must point at a node, and an
% `example` is not one -- linkage does not count it a statement environment and the dependency
% graph does not carry it.
%
% One node per family, not a definition/proposition pair. A pair would occupy 8.1-8.6 and push
% Proposition 8.4 to 8.7, breaking the numbering agreement with the draft for the sake of a
% distinction the content does not need: each of these states "this F is admissible and its
% kernels are these", which is one claim. Same reasoning as \ref{prop:bernstein-toolbox}.
%
% Two [A] nodes here, the first since chapter 2: \ref{prop:bessel-family} on ledger A8
% (Halgreen) and \ref{prop:moment-criterion} on A7 (Sato). Both say more than their citation carries,
% so both Assignment clauses divide the statement explicitly.

Throughout this chapter we work in the canonical gauge, so that
$\phi_x = \Lap^{-1}[e^{-F(x s)}]$, and $T_x$ denotes a random variable with density $\phi_x$.

% draft: Example 8.1
\begin{proposition}[the extremal stable family; the 2005 kernels]\label{prop:stable-family}
\uses{lem:selfdecomposable-exponents,cor:semigroup-case,prop:moment-criterion}
For $0 < \alpha < 1$ the exponent $F(s) = s^\alpha$ is of the form (7.1), with Lévy density
factor
\[
  k(t) = \frac{\alpha}{\Gamma(1-\alpha)}\, t^{-\alpha},
\]
which is nonincreasing. Every moment of $T_x$ is infinite, so the delay cannot be measured by
the mean; measured by the mode it scales as $x^{1/\alpha}$.
\statusT
\end{proposition}

\begin{proof}
$s^\alpha = \frac{\alpha}{\Gamma(1-\alpha)}\int_0^\infty (1 - e^{-st})\,t^{-\alpha-1}\,dt$ is
the standard integral representation, which is (7.1) with $b_0 = 0$ and
$k(t)/t = \alpha t^{-\alpha-1}/\Gamma(1-\alpha)$; $k$ is a negative power, hence
nonincreasing. Since $\int_1^\infty t^{n-1}k(t)\,dt = \infty$ already for $n = 1$, all moments
are infinite by \ref{prop:moment-criterion}. The mode scaling is the self-similarity
$T_x \stackrel{d}{=} x^{1/\alpha} T_1$ of the $\alpha$-stable subordinator, which is the case
$\chi(x) = x^{1/\alpha}$ of \ref{cor:semigroup-case}.
\end{proof}

% draft: Example 8.2
\begin{proposition}[the Gamma family]\label{prop:gamma-family}
\uses{lem:selfdecomposable-exponents,thm:main-characterization,rem:sato-reading}
For $\gamma > 0$ the exponent $F(s) = \gamma\,\log(1 + s)$ is of the form (7.1) with
$k(t) = \gamma\, e^{-t}$, which is nonincreasing. Its kernels are
\[
  e^{-F(xs)} = (1 + x s)^{-\gamma}, \qquad
  \phi_x(t) = \frac{t^{\gamma - 1} e^{-t/x}}{\Gamma(\gamma)\, x^{\gamma}},
\]
the Gamma density of fixed shape $\gamma$ and scale $x$. All moments are finite, with
$\Ex\,T_x = \gamma x$ and $\operatorname{Var} T_x = \gamma x^2$. The increment measures
$\mu_{a,b}$ have Laplace transform $\bigl(\tfrac{1 + as}{1 + bs}\bigr)^{\gamma}$ and are the
increments of the Gamma-type Sato subordinator.
\statusT
\end{proposition}

\begin{proof}
The Frullani integral gives
$\log(1+s) = \int_0^\infty (1 - e^{-st})\,e^{-t}\,t^{-1}\,dt$, so $F$ has the form (7.1) with
$b_0 = 0$ and $k(t) = \gamma e^{-t}$, nonincreasing; hence $F$ is admissible by
\ref{lem:selfdecomposable-exponents}. Then $e^{-F(xs)} = (1+xs)^{-\gamma}$, which is the
Laplace transform of the stated Gamma density. The moments are those of a Gamma law, and all
are finite because $\int_1^\infty t^{n-1}e^{-t}dt < \infty$ for every $n$. The increment
transform is the ratio $e^{-F(bs)}/e^{-F(as)}$, and \ref{rem:sato-reading} identifies the
family as the marginals of a Sato subordinator.
\end{proof}

% draft: Example 8.3
\begin{proposition}[the Bessel-K family; preview of the memory line]\label{prop:bessel-family}
\uses{lem:selfdecomposable-exponents,thm:main-characterization,prop:moment-criterion}
For $a > 0$ let $T_1 \stackrel{d}{=} 1/(2\gamma_a)$ be the scaled inverse-gamma law, where
$\gamma_a$ has the Gamma density of shape $a$, and let $F$ be its Laplace exponent. Then
\[
  e^{-F(s)} = \frac{2^{1-a}}{\Gamma(a)}\,\bigl(\sqrt{2s}\bigr)^{a} K_a\bigl(\sqrt{2s}\bigr),
\]
and $T_1$ is self-decomposable, so $F$ is of the form (7.1) and the family is admissible. In
the parabolic gauge $\tilde x = x^2/2$ these are the boundary-hitting kernels of the
Bessel-type media $\tfrac12\partial_x^2 + \tfrac{\beta}{x}\partial_x$ with $a = \tfrac12 -
\beta$, and the moments satisfy $\Ex\,T_x^n < \infty$ if and only if $a > n$. This family is
the subject of the memory-line development, Theorems 10.4, 11.6 and 12.5.
\statusA\quad\ledger{A8}\quad\emph{(self-decomposability of the generalized inverse Gaussian
laws; \sscite{halgreen1979self}. A cited interface, not reproven here.
\textbf{Assignment.} \ledger{A8} carries exactly one clause: that $T_1$ is self-decomposable,
which is what makes $F$ admissible and puts the family inside
\ref{thm:main-characterization}. Everything else is \textbf{[T]} and proved below --- the
evaluation of the transform in terms of $K_a$ is a classical integral representation, and the
moment criterion is \ref{prop:moment-criterion} applied to this $k$. The parabolic-gauge reading is a
forward reference to Chapter~11 and is not asserted here. Note also that Halgreen proves the
stronger statement that the law is a generalized Gamma convolution; self-decomposability
follows through $\GGC \subset \SDclass$, which is a clause of \ledger{A16} --- so A8 and A16
rest on the same tower and should be reviewed together.)}
\end{proposition}

\begin{proof}
Write $\gamma_a$ for a Gamma$(a,1)$ variable and $T_1 = 1/(2\gamma_a)$. Substituting
$u = 1/(2t)$ in the Gamma density gives the density of $T_1$,
\[
  \phi_1(t) = \frac{t^{-a-1}\,e^{-1/(2t)}}{2^{a}\,\Gamma(a)}, \qquad t > 0 .
\]
The classical integral representation
$\int_0^\infty t^{\nu-1}\exp(-A/t - Bt)\,dt = 2\,(A/B)^{\nu/2}K_\nu\bigl(2\sqrt{AB}\bigr)$,
taken at $\nu = -a$, $A = \tfrac12$, $B = s$, together with $K_{-a} = K_a$, gives
\[
  \Ex\,e^{-sT_1}
  = \frac{1}{2^{a}\Gamma(a)}\int_0^\infty t^{-a-1}e^{-st - 1/(2t)}\,dt
  = \frac{2\,(2s)^{a/2}\,K_a\bigl(\sqrt{2s}\bigr)}{2^{a}\Gamma(a)}
  = \frac{2^{1-a}}{\Gamma(a)}\bigl(\sqrt{2s}\bigr)^{a}K_a\bigl(\sqrt{2s}\bigr).
\]
Self-decomposability of $T_1$ is \ledger{A8}; by
\ref{lem:selfdecomposable-exponents} its exponent is therefore of the form (7.1), and
\ref{thm:main-characterization} makes the family admissible. Finally, the Lévy density factor
of this $F$ satisfies $k(t) \sim a$ as $t \to 0$ and decays like $t^{-a}$ at infinity, so
$\int_1^\infty t^{n-1}k(t)\,dt < \infty$ exactly when $n < a$; \ref{prop:moment-criterion} converts
that into $\Ex\,T_x^n < \infty$ iff $a > n$.
\end{proof}

% draft: Proposition 8.4
\begin{proposition}[the mean delay and the influence curve]\label{prop:moments}
\uses{lem:selfdecomposable-exponents,thm:main-characterization}
$\Ex\,T_x = x\,F'(\zp) = x\bigl(b_0 + \int_0^\infty k(t)\,dt\bigr)$, finite if and only if $k$
is integrable at infinity. When the mean is finite, the influence curve of
\sscite{fagerstrom2005temporal}, \S7, may be defined by the mean delay, and it is exactly
linear in the canonical gauge:
\[
  t_0 - t_m = F'(\zp)\; x .
\]
\statusT\quad\emph{(monotone convergence in (7.1) and nothing else --- no citation is needed and
none is made. The higher-moment criterion, which \emph{does} need one, was split off as
\ref{prop:moment-criterion}; before the split this node reported that entry's cost for a result
that does not incur it. What is \emph{not} claimed here is anything about the mode, which is the
delay measure the stable family forces one to use --- see \ref{prop:stable-family}.)}
\end{proposition}

\begin{proof}
Differentiating the transform at the origin,
$\Ex\,T_x = -\partial_s e^{-F(xs)}\big|_{s=\zp} = x F'(\zp)$, and monotone convergence in
(7.1) gives $F'(\zp) = b_0 + \int_0^\infty k(t)\,dt$, finite exactly when $k$ is integrable at
infinity. The influence curve is then linear in $x$ because $\Ex\,T_x$ is.
\end{proof}

% draft: Remark 8.5
\begin{remark}[heavy tails are the semigroup's doing]\label{rem:heavy-tails}
\uses{thm:main-characterization,cor:semigroup-case,prop:moment-criterion,prop:gamma-family}
The heavy tails of the 2005 kernels are an artifact of the semigroup axiom, not of causality
together with covariance. Within \ref{thm:main-characterization} the semigroup sub-case
(\ref{cor:semigroup-case}) is precisely the subset with $k$ a pure power --- the only $k$
compatible with homogeneity --- and every non-power nonincreasing $k$, for instance
$k(t) = \gamma e^{-t}$, yields a causal, continuously scale-covariant family with finite
moments of all orders.
\end{remark}

% additive: not in the draft. The higher-moment half of Proposition 8.4, split off so that the
% mean-delay identity -- which needs no citation -- does not report ledger A7's cost. Numbered
% 8.6, a number the draft does not use, its Section 8 ending at 8.5.
\begin{proposition}[the moment criterion]\label{prop:moment-criterion}
\uses{prop:moments,lem:selfdecomposable-exponents,thm:main-characterization}
$\Ex\,T_x^{\,n} < \infty$ if and only if $\int_1^\infty t^{\,n-1} k(t)\,dt < \infty$.
\statusA\quad\ledger{A7}\quad\emph{(the moment criterion for infinitely divisible laws;
\sscite{sato1999levy}, Thm.~25.3. A cited interface, not reproven here.
\textbf{Assignment.} \ledger{A7} carries the equivalence itself and nothing else --- that a
moment of the law is finite exactly when the corresponding integral against the Lévy measure
over $\{t > 1\}$ is finite. Translating that integral into the stated one is \textbf{[T]}: the
Lévy measure of $T_x$ is $\nu(dt) = k(t)t^{-1}dt$ up to the scaling by $x$, so
$\int_{t>1}t^n\,\nu(dt) = \int_1^\infty t^{\,n-1}k(t)\,dt$. The $n = 1$ case is \emph{not}
carried by the ledger --- it is \ref{prop:moments}, proved by monotone convergence --- which is
the reason for the split.)}
\end{proposition}

\begin{proof}
The Lévy measure of $T_x$ is $\nu(dt) = k(t)t^{-1}dt$ up to the scaling by $x$, so
$\int_{t>1} t^n\,\nu(dt) = \int_1^\infty t^{\,n-1}k(t)\,dt$; \ledger{A7} equates finiteness of
that integral with finiteness of $\Ex\,T_x^n$.
\end{proof}


\chapter{The signaling form: Theorem 4$'$}
% !TEX root = ../print.tex
% Draft \S11 --- The signaling form (Theorem 4')
%
% Source: draft/hemigroup-causal-scale-space-kernels.md, lines 465-566.
% This file is chapter 11 --- The signaling form (Theorem 4') of the blueprint; the shared counter makes every number here
% agree with the draft's own (see content.tex).
%
% Planned nodes: def:standing-hypothesis, lem:mellin-data, def:inversion-operator, lem:symbol-uniqueness, lem:memory-fractional-integrals, thm:signaling-form [T]; Examples 11.8-11.10, Remarks 11.6/11.7/11.11
%
% No nodes yet. Each statement carries \label{kind:name}, \lean{...} with
% \leanok/\notready, \uses{...}, \notes{slug}, a \statusT or \statusA line, and a
% "% draft: <kind> N.M" provenance comment; an \statusA node must also declare what its
% citation carries (\textbf{Assignment.}).


\chapter{The locality theorem: Theorem 5$'$}
% !TEX root = ../print.tex
% Draft \S12 --- The locality theorem (Theorem 5')
%
% Source: draft/hemigroup-causal-scale-space-kernels.md, lines 567-638.
% This file is chapter 12 --- The locality theorem (Theorem 5') of the blueprint; the shared counter makes every number here
% agree with the draft's own (see content.tex).
%
% Planned nodes: def:locality-pmp, lem:local-polynomial-symbol, lem:moment-recursion, lem:log-convexity, thm:locality [T], prop:local-ladder; Remarks 12.7/12.8
%
% No nodes yet. Each statement carries \label{kind:name}, \lean{...} with
% \leanok/\notready, \uses{...}, \notes{slug}, a \statusT or \statusA line, and a
% "% draft: <kind> N.M" provenance comment; an \statusA node must also declare what its
% citation carries (\textbf{Assignment.}).


\chapter{Implementation}
% !TEX root = ../print.tex
% Draft \S13 --- Implementation
%
% Source: draft/hemigroup-causal-scale-space-kernels.md, lines 705-804, as revised by the
% paper's review-response batches (R5's graded substrate statuses, the drift caveat).
% Chapter 13 of the blueprint; the shared counter puts every number on the draft's own.
%
% ADDED 2026-09-01 on the author's instruction, REVERSING the original scope decision
% (content.tex's header): the implementation chapter was excluded early to keep the
% blueprint's size under control; the author now wants the blueprint complete with
% respect to the mathematical material, with formalisation a separate decision. So the
% nodes below are blueprint statements first and formalisation targets not at all for
% now: none carries a \lean tag, and none is scheduled in PLAN-chapters-8-12.md.
%
% One ledger entry lands here: A21 (Thorin's characterization of the GGC class), the
% analytic substance of \ref{prop:thorin-subclass}; its (1) <=> (2) leg rides on A1's
% general-measure Bernstein--Widder form. Anchors image-verified 2026-09-01.

This chapter records how the scale spaces of \ref{thm:main-characterization} are computed.
The structural fact that organizes everything: the scale direction never needs to be
discretized approximately.

% draft: Proposition 13.1
\begin{proposition}[exactness at the knots]\label{prop:knot-exactness}
\uses{lem:convolution-representation}
Let $0 = x_0 < x_1 < \cdots < x_M$ be any finite set of scale knots and define the discrete
cascade
\[
  u_0 := f, \qquad u_{k+1} := \mu_{x_k,\,x_{k+1}} * u_k .
\]
Then $u_k = \Phi_{0, x_k} f$ \emph{exactly}, for every knot set; and refining the knot set
changes nothing at the old knots.
\statusT\quad\emph{(One line of hemigroup algebra; the content is that the cascade is not an
approximation scheme for the family but the family itself, sampled --- the only
discretization error in a digital implementation is the time discretization of the
increment filters.)}
\end{proposition}

\begin{proof}
Immediate induction from the hemigroup (A6): $\mu_{0,x_{k+1}} = \mu_{x_k, x_{k+1}} * \mu_{0,x_k}$.
\end{proof}

% draft: Proposition 13.2
\begin{proposition}[the Gamma cascade]\label{prop:gamma-cascade}
\uses{prop:gamma-kernels,lem:convolution-representation}
For the Gamma family with integer $\gamma$, the increment filter is rational:
\[
  \widehat{\mu}_{x_k, x_{k+1}}(s) \;=\; \Bigl(\frac{1 + x_k\, s}{1 + x_{k+1}\, s}\Bigr)^{\gamma},
\]
a cascade of $\gamma$ \emph{identical} first-order pole--zero sections (pole at $-1/x_{k+1}$,
zero at $-1/x_k$). Each section admits the one-state realization, free of input differentiation,
\[
  \dot w = \frac{y - w}{x_{k+1}}, \qquad
  v \;=\; \frac{x_k}{x_{k+1}}\, y \;+\; \Bigl(1 - \frac{x_k}{x_{k+1}}\Bigr)\, w,
\]
($y$ input, $v$ output), whose transfer is $(1 + x_k s)/(1 + x_{k+1} s)$. The output is a
\emph{convex combination} of input and state.
\statusT\quad\emph{(Transform algebra; the convex combination is what makes positivity and
unit mass manifest at the realization level. In discrete time, with sample interval
$\Delta$ and sampled-and-held input, each section admits the exact exponential-integrator
update $w_{n+1} = e^{-\Delta/x_{k+1}} w_n + (1 - e^{-\Delta/x_{k+1}}) y_n$ ---
unconditionally stable and free of overshoot; the cascade is not exact end-to-end, the
input to the next section being the previous section's response rather than a held
signal, so time discretization is a genuine approximation, while positivity, unit mass
and stability are preserved exactly by the convex update.)}
\end{proposition}

\begin{proof}
The increment transfer is the ratio of the endpoint transforms: by the hemigroup (A6),
$\widehat{\mu}_{0,x_{k+1}} = \widehat{\mu}_{x_k,x_{k+1}}\,\widehat{\mu}_{0,x_k}$, and
$\widehat{\mu}_{0,x}(s) = (1 + x s)^{-\gamma}$ is nonvanishing for $\operatorname{Re} s \ge 0$,
so dividing gives the display. The transfer identity of the section:
$v = \tfrac{x_k}{x_{k+1}} y + \bigl(1 - \tfrac{x_k}{x_{k+1}}\bigr)\tfrac{y}{1 + x_{k+1}s}
= y\,\tfrac{1 + x_k s}{1 + x_{k+1} s}$; the convex-combination clause is
$0 < x_k/x_{k+1} < 1$.
\end{proof}

% draft: Remark 13.3; the paper's copy carries the sampled-system precisions added on
% review (time constants tied to the sample interval; re-instantiation, not relabelling).
\begin{remark}[covariance in practice]\label{rem:covariance-practice}
\uses{prop:gamma-cascade}
With geometric knots $x_k = q^k x_0$ the \emph{continuum} cascade is carried to itself by
every rescaling $\sigma \in q^{\ZZ}$, as a shift of the ladder; in the sampled system the
time constants are tied to the sample interval $\Delta$, so the shift is exact on the knot
set but not on the sampled filters. Because the underlying family exists at \emph{every}
scale, an arbitrary $\sigma$ is handled by instantiating the same family at the shifted
knots --- a new filter bank with new time constants and re-initialized state, not a
relabelling of the running one; the operation is principled exactly because the kernels
between knots belong to the family. This is the operational content of continuous
covariance, and the point of the pole--zero structure: a pure-pole cascade (truncated
exponentials with distinct time constants), the finite form of the time-causal limit-kernel
construction of \sscite{lindeberg2016time}, also marches a hemigroup, but its kernels are
hypoexponential laws whose family is \emph{not} dilation-closed, so only the discrete
covariance of the ladder itself survives. The zeros are what the hemigroup increments of the
Sato--Gamma process add, and they are exactly what buys the continuum.
\end{remark}

% draft: Remark 13.4; clause (iii) as revised by R5 (the state-count claim deferred to the
% heuristic paragraph after \ref{prop:thorin-subclass}).
\begin{remark}[implementing the local corner]\label{rem:local-implementation}
\uses{prop:bessel-family,prop:stable-family,thm:locality,prop:thorin-subclass,rem:gauge-freedom}
The Bessel/inverse-gamma increments have transcendental transforms (ratios of Bessel-K
profiles), hence \emph{no} finite-dimensional time-recursive realization; the same holds for
the stable family ($e^{-(x_{k+1}^{\,\alpha} - x_k^{\,\alpha})\,s^\alpha}$ in the homogeneous
gauge). The implementation options are: (i) direct convolution against tabulated increment
kernels per knot --- burdened by the polynomial tails $t^{-1-a}$, so truncation error decays
only polynomially in the window length; (ii) marching the memory-line PDE, in the parabolic
gauge $\xi$ of \ref{rem:gauge-freedom},
$\partial_t u = \tfrac12 u_{\xi\xi} + \tfrac{\beta}{\xi}u_\xi$ in time (implicit in $\xi$, with
care at the degenerate boundary), which reintroduces scale-direction discretization error; or
(iii) rational approximation of the irrational symbol, in the manner standard for
fractional-order elements in filter and control design (\sscite{oustaloup2000frequency});
\ref{prop:thorin-subclass} keeps that route inside the axiom class, and its state-count
economy is a heuristic, with no error bound claimed. All three are workable; none matches
the cascade of \ref{prop:gamma-cascade} in simplicity or exactness. This is the practical
face of \ref{thm:locality}: locality of the medium is bought at the price of heavy-tailed,
non-rational kernels.
\end{remark}

% draft: Remark 13.5
\begin{remark}[finite differences on the memory line]\label{rem:finite-differences}
\uses{prop:bessel-family,prop:knot-exactness,rem:gauge-freedom}
The explicit finite-difference scheme of \sscite{fagerstrom2005temporal}, \S\S6--7, for the
heat signaling equation extends verbatim to the Bessel family. Working in the parabolic gauge
$\xi$ of \ref{rem:gauge-freedom}, on the logarithmic grid $\eta = \log\xi$ the operator
is constant-coefficient up to a rate factor,
\[
  \tfrac12\,\partial_\xi^2 + \frac{\beta}{\xi}\,\partial_\xi
  \;=\; e^{-2\eta}\Bigl[\tfrac12\,\partial_\eta^2 + \bigl(\beta - \tfrac12\bigr)\,\partial_\eta\Bigr],
\]
so the stencil is uniform along the scale ladder and only the update rate varies with
position --- the discrete form of the Lamperti time change (\sscite{lamperti1972semi}). With
the first-order term upwinded, the explicit update matrix is sub-stochastic, so positivity
and non-amplification are automatic, and the discrete scheme is itself a hemigroup of
sub-probability kernels: a random-walk medium whose absorption-time laws are discrete
analogues of the inverse-gamma kernels. It is recorded for continuity with the 2005 paper but
not recommended: the explicit step is constrained by the finest cell, to steps $\Delta t$
of order $\xi_{\min}^2\,\Delta\eta^2$, while the coarse-scale dynamics live on
horizons $\sim \xi_{\max}^2$, so the cost of covering the ladder grows like
$(\xi_{\max}/\xi_{\min})^2$ (an implicit march trades this for per-sample band solves plus
additional discretization error);
the finite-moment regime $a > 1$, i.e.\ $\beta < -\tfrac12$ --- the very reason to prefer the
Bessel family --- is advection-dominated near the singular boundary and numerically the least
pleasant; and, decisively, any marching of the signaling equation approximates the scale
direction, which \ref{prop:knot-exactness} shows can be sampled exactly. The principled
numerics for the local corner is instead the following.
\end{remark}

% draft: Proposition 13.6, as revised by R5 (the graded substrate statuses live in the
% paper's discussion after the proof; the statement itself is unchanged in substance,
% with the drift rider softened to "a pure time shift").
\begin{proposition}[the Thorin subclass; gamma cascades approximate the local corner]\label{prop:thorin-subclass}
\uses{thm:main-characterization,lem:memory-kernel,lem:selfdecomposable-exponents,prop:laplace-continuity,prop:pair-regularity,prop:gamma-cascade}
For an admissible $F$ (\ref{thm:main-characterization}) the following are equivalent:
\begin{enumerate}
\item the memory kernel $k$ is completely monotone --- the classical Sonine case of
\ref{prop:pair-regularity}(2);
\item $F$ has a \emph{Thorin representation}
\[
  F(s) \;=\; b_0\,s \;+\; \int_{(0,\infty)} \log\Bigl(1 + \frac{s}{\tau}\Bigr)\,U(d\tau),
  \qquad U \ge 0,\ \
  \int_{(0,1]} \log\tfrac{1}{\tau}\,U(d\tau) + \int_{(1,\infty)} \frac{U(d\tau)}{\tau} < \infty;
\]
\item the kernels are generalized gamma convolutions (GGC).
\end{enumerate}
In that case, for any atomic $U_N = \sum_{i=1}^N w_i\,\delta_{\tau_i}$ such that
$F_N(s) := b_0 s + \sum_i w_i \log(1 + s/\tau_i)$ converges to $F$ pointwise --- as it does
for the standard quadratures of $U$, the weights carrying the two integrability conditions
--- the exponents $F_N$ are \emph{themselves admissible}, finite multi-Gamma members of
\ref{thm:main-characterization}, and the kernel families converge weakly at every
scale; positivity and unit mass hold \emph{exactly at every truncation level}, the
approximation staying inside the axiom class rather than merely near it. For integer weights
the increments over any knot ladder are rational of order $\sum_i w_i$ up to the pure delay
$e^{-b_0(x_{k+1}-x_k)s}$ contributed by the drift --- itself a pure time shift --- realized
by the cascade of \ref{prop:gamma-cascade} applied per
Thorin node, with time constants $x/\tau_i$.
\statusA\quad\ledger{A1}\quad\ledger{A21}\quad\emph{(Thorin's characterization of the GGC
class, plus Bernstein--Widder.
\textbf{Assignment.} \ledger{A21} carries the equivalence (2) $\Leftrightarrow$ (3): a law
on $[0,\infty)$ is GGC iff it is infinitely divisible with L\'evy density $\ell$ such that
$y\,\ell(y)$ is completely monotone (\sscite{bondesson1992generalized}, Thm.~3.1.1, p.~30
--- clause (1) verbatim, $y\,\ell(y)$ being this development's $k$), and, exponent-side,
the Thorin--Bernstein form with its unique representing measure and exactly the displayed
integrability condition (\sscite{schilling2012bernstein}, Thm.~8.2, p.~109, (i)
$\Leftrightarrow$ (iii)). \ledger{A1}, in its general-measure form, carries the step $k$
completely monotone $\Rightarrow$ $k(t) = \int e^{-\tau t}\,U(d\tau)$ inside (1)
$\Rightarrow$ (2). Everything else is \textbf{[T]} in prose: the integration of $F'$
against \ref{lem:memory-kernel}, the three-regime finiteness bookkeeping, the
admissibility of the truncations through \ref{lem:selfdecomposable-exponents} and
\ref{thm:main-characterization}, and the weak convergence through
\ref{prop:laplace-continuity}. The graded evidentiary statuses of the substrate claim ---
in-class approximation proved at any positive weights, exact rational realization proved
at integer weights, real-weight fitting a heuristic --- are the paper's discussion, not
part of this statement.)}
\end{proposition}

\begin{proof}
(1) $\Leftrightarrow$ (2): by Bernstein--Widder (\ledger{A1}; in the working source,
\sscite{schilling2012bernstein}, Thm.~1.4, p.~3), $k$ completely monotone means
$k(t) = \int e^{-\tau t}\,U(d\tau)$ for a
positive measure $U$ on $[0,\infty)$, finite on compacts since the integral converges; an
atom of $U$ at $\tau = 0$ would put a constant component into $k$, making
$\int_1^\infty k(t)\,t^{-1}dt$ --- hence $F$ --- infinite, so $U$ lives on $(0,\infty)$. By
\ref{lem:memory-kernel},
$F'(s) = b_0 + \int U(d\tau)/(s + \tau)$, and integrating in $s$ from $0$ gives the Thorin
form, the integrability condition being finiteness of $F(1)$ --- three regimes:
$\log(1+1/\tau) \sim \log(1/\tau)$ as $\tau \downarrow 0$, $\sim 1/\tau$ as
$\tau \to \infty$, and bounded on the compact middle, where $U$ is finite; conversely
differentiate.
(2) $\Leftrightarrow$ (3) is Thorin's characterization of the GGC class
(\ledger{A21}; \sscite{thorin1977pareto}, \sscite{thorin1977lognormal}). In the L\'evy-data
form --- GGC iff the L\'evy density $\ell$ has $y\,\ell(y)$ completely monotone, which is
clause (1) verbatim, $y\,\ell(y)$ being this development's $k$ --- it is
\sscite{bondesson1992generalized}, Thm.~3.1.1, p.~30; the exponent-side form is the
Thorin--Bernstein class of \sscite{schilling2012bernstein}, Thm.~8.2, p.~109.
Admissibility of $F_N$: $sF_N'(s) = b_0 s + \sum_i w_i\, s/(s + \tau_i)$ is a Bernstein
function, so \ref{lem:selfdecomposable-exponents} applies and $F_N$ is of the form
(7.1); $F_N(\zp) = 0$, the integrability conditions are trivial for a finite sum, and
\ref{thm:main-characterization} makes the family admissible. Pointwise
convergence of $F_N$ gives convergence of the transforms $e^{-F_N(xs)}$, hence weak
convergence of the kernels by the continuity theorem
(\ref{prop:laplace-continuity}).
\end{proof}

% draft: Example 13.7. Condensed from the paper's ex:tail-contrast: the figure panels and
% the closing pointer to the paper's end-to-end machine run (ex:numerical-experiment,
% paper-only) are omitted; the two window-length tables, the quantitative content, are
% kept verbatim.
\begin{example}[numerical illustration: the tail contrast]\label{ex:tail-contrast}
\uses{prop:gamma-kernels,prop:bessel-family,prop:stable-family}
Compare the kernels of the three corners at matched delay: Gamma kernels with
$\gamma \in \{1, 2, 4\}$ and scale $x = 1/\gamma$ (mean delay $1$); the Bessel/inverse-gamma
kernel with $a = 2$, normalized to mean delay $1$;
and the $\tfrac12$-stable kernel of \sscite{fagerstrom2005temporal} normalized to median
delay $1$, its mean being infinite. At matched mean the Gamma kernels are the most
concentrated, sharpening with $\gamma$ (mode $(\gamma - 1)/\gamma \to$ mean), while the
inverse-gamma kernel buys its early sharp peak by exporting mass into the tail, and the
stable kernel does so to the point of losing the mean altogether; the Gamma tails fall
exponentially, the inverse-gamma and stable tails follow the powers $t^{-1-a}$ and
$t^{-3/2}$ of \ref{rem:trilemma}. The practical consequence is quantified by the
window length $T_p$ capturing a fraction $p$ of the kernel mass ($\int_0^{T_p}\phi = p$),
the quantity governing the truncation error of route (i) in \ref{rem:local-implementation}.
At matched delay $1$, the pairs $(T_{0.99},\, T_{0.999})$ are:
\begin{itemize}
\item $\Gamma(\gamma{=}1)$: $(4.61,\ 6.91)$; $\Gamma(\gamma{=}2)$: $(3.32,\ 4.62)$;
  $\Gamma(\gamma{=}4)$: $(2.51,\ 3.27)$;
\item inverse-gamma, $a = 2$: $(6.73,\ 22.0)$;
\item $\tfrac12$-stable: $(2.90 \times 10^{3},\ 2.90 \times 10^{5})$.
\end{itemize}
For the Gamma family each extra decade of captured mass costs a fixed additive window
increment (exponential tails); for the inverse-gamma with $a = 2$ it costs a factor
$10^{1/a}$ (about $3.2$); for the $\tfrac12$-stable, of index $\alpha = \tfrac12$, a factor
$10^{1/\alpha} = 100$ --- capturing
$99.9\%$ of its mass requires a window five orders of magnitude longer than its median
delay. Since the stable mean is infinite, the delay proxy of the influence-curve discussion
of \sscite{fagerstrom2005temporal}, \S8, is the \emph{mode}; repeating the comparison at
matched mode delay $1$ ($x = 1/(\gamma-1)$ for the Gamma kernels, the inverse-gamma and
stable normalized likewise) gives, as triples $(\text{mean},\, T_{0.99},\, T_{0.999})$:
\begin{itemize}
\item $\Gamma(\gamma{=}2)$: $(2,\ 6.64,\ 9.23)$; $\Gamma(\gamma{=}4)$: $(4/3,\ 3.35,\ 4.35)$;
\item inverse-gamma, $a = 2$: $(3,\ 20.2,\ 66.1)$;
\item $\tfrac12$-stable: $(\infty,\ 1.91 \times 10^{4},\ 1.91 \times 10^{6})$
\end{itemize}
--- at matched mode, capturing $99.9\%$ of the $\tfrac12$-stable kernel requires a window
six orders of magnitude beyond the peak. This is the quantitative content of the statement
that the 2005 kernels are too heavy-tailed for practical measurement, and of its resolution
within the hemigroup class.
\end{example}


\chapter{The temporal N-jet: the memory embodies its derivatives}
% !TEX root = ../print.tex
% Draft \S14 --- The temporal N-jet: the memory embodies its derivatives
%
% Source: draft/hemigroup-causal-scale-space-kernels.md, lines 806-848, as revised by the
% paper (the jet-taps architecture discussion, the Rodrigues form of rem:laguerre --
% the draft's orthogonality claim was corrected on numerical verification, 2026-08).
% Draft §14; renders as Chapter 12 since the 2026-09-01 reorder to the article's order
% (content.tex has the rationale). Historical comments below that cite chapter or node
% numbers ("Numbered 14.M") use the retired draft-aligned numbering; the `% draft:`
% provenance comments are unaffected, and labels and \uses edges never carried numbers.
%
% ADDED 2026-09-01 together with chapter 13, on the author's instruction (see that file's
% header): blueprint completeness first, formalisation a separate decision. No node below
% carries a \lean tag. Everything rests on chapter 11's Mellin apparatus and chapter 13's
% cascade; no new ledger entry lands here.
%
% Two statements were lightly restructured IN THE PAPER for verbatim sharing (2026-09-01):
% thm:embodied-jet's "as in \S9" became "as at order one", and prop:gamma-jets' trailing
% architecture/figure/strip-extension discussion moved out of the statement into body
% text, on the prop:bessel-family precedent (2026-08-23). The paper's machine-run example
% (ex:numerical-experiment) stays paper-only: it is a run, not a statement.

The jet vocabulary iterates the operator core: for $n \ge 1$ let
\[
  \core^n \;:=\; \bigl\{\, f \in \Xz : \text{there are } f^{(1)}, \dots, f^{(n)} \in \Xz
  \text{ with } f^{(j)} = \textstyle\int_0^{\,\cdot} f^{(j+1)},\ 0 \le j < n \,\bigr\},
\]
so that $\core^1 = \core$: each derivative up to order $n$ is again a causal $L^1$
function that is the primitive of the next, and the vanishing boundary values are exactly
what turns $I^{z}f^{(n)}$ into $I^{z-n}f$ below.

% draft: Definition 14.1
\begin{definition}\label{def:jet}
\uses{prop:canonical-gauge}
The \emph{temporal $N$-jet} of the scale space at $(t, x)$ is
$\{\partial_t^n u(t, x)\}_{0 \le n \le N}$, and the \emph{scale-normalized jet} is
$\{x^n \partial_t^n u(t, x)\}_{n \le N}$, which in the canonical gauge transforms under the
joint dilation $(t, x) \mapsto (\sigma t, \sigma x)$ like the field itself.
\statusT\quad\emph{(A definition; its one claim --- the transformation law of the
scale-normalized jet --- is the chain rule under the joint dilation, read off the
canonical-gauge covariance of \ref{prop:canonical-gauge}, each $\partial_t$ contributing a
factor $\sigma^{-1}$ that the weight $x^n$ restores.)}
\end{definition}

% draft: Theorem 14.2
\begin{theorem}[the embodied jet]\label{thm:embodied-jet}
\uses{def:jet,def:standing-hypothesis,lem:memory-fractional-integrals,def:inversion-operator,rem:poles}
Assume (H), canonical gauge, $f \in \core^n$, and let $u(\cdot, x) = \Phi_{0,x}f$. Then:
\begin{enumerate}
\item For every $t > 0$ and $1 < \operatorname{Re} z < z_*$,
\[
  \widetilde{\partial_t^n u(t, \cdot)}(z) \;=\; \tilde H(z)\,\bigl(I^{z} f^{(n)}\bigr)(t)
  \;=\; \tilde H(z)\,\bigl(I^{\,z - n} f\bigr)(t),
\]
the second equality holding as a Riemann--Liouville integral for $\operatorname{Re} z > n$ and
by analytic continuation --- a fractional derivative of $f$ --- below.
\item Consequently, whenever $z_* > n + 1$, the $n$-th temporal derivative is the $n$-th power
of the inversion operator applied to the \emph{present} memory state: in the Mellin sense on
$n + 1 < \operatorname{Re} z < z_*$,
\[
  \partial_t^n u(t, \cdot) \;=\; A^n\, u(t, \cdot), \qquad
  A^n = \frac{1}{x^n}\,B^{(n)}(\theta), \qquad
  B^{(n)}(-z) := \frac{\tilde H(z + n)}{\tilde H(z)} = \prod_{j=0}^{n-1} B(-z - j)
\]
for $0 < \operatorname{Re} z < z_* - n$ --- the identity evaluating the symbol at the shifted
argument, as at order one.
\end{enumerate}
In particular the temporal jet at time $t$, up to every order $n < z_* - 1$, is an
instantaneous linear functional of the single profile $x \mapsto u(t, x)$; no record of past
states is required. The identity is Mellin-level, with the pointwise reading subject to the
same pole discipline as order one (\ref{rem:poles}).
\statusT\quad\emph{(The embodiment principle as theorem content: differentiation is an
instantaneous spatial readout of the present memory state, no second memory recording the
field's history. The proof is an induction through \ref{lem:memory-fractional-integrals}'s
derivative clause plus a telescoping of \ref{def:inversion-operator}'s shift relation ---
chapter 9's apparatus, nothing new taken on trust.)}
\end{theorem}

\begin{proof}
(1) By induction on the order: for $j < n$, $f^{(j)}$ is the primitive of
$f^{(j+1)} \in \Xz$, so \ref{lem:memory-fractional-integrals}'s derivative clause
applies at each step --- $\partial_t\,\Ex[f^{(j)}(\cdot - xT_1)] = \Ex[f^{(j+1)}(\cdot - xT_1)]$
as $\Xz$-derivatives, the boundary conditions of $\core^n$ licensing each step --- and $n$
steps give $\partial_t^n u(t, x) = \Ex\bigl[f^{(n)}(t - x T_1)\bigr]$. The lemma applied to
the causal $f^{(n)} \in \Xz$ then gives the
first equality; the second is $I^{z} f^{(n)} = I^{\,z-n} I^{\,n} f^{(n)} = I^{\,z-n} f$,
$I^{\,n}f^{(n)} = f$ being the primitive structure of $\core^n$ again. (2) Telescoping the shift relation of
\ref{def:inversion-operator} $n$ times,
$\Mel[A^n g](z) = \prod_{j=1}^{n} B(j - z)\,\tilde g(z - n)
= \tfrac{\tilde H(z)}{\tilde H(z - n)}\,\tilde g(z - n)$; insert
$\tilde g = \widetilde{u(t,\cdot)}$ from \ref{lem:memory-fractional-integrals} and
compare with (1).
\end{proof}

% draft: Remark 14.3
\begin{remark}[differentiation is re-indexing]\label{rem:reindexing}
\uses{thm:embodied-jet,rem:embodiment}
Combined with \ref{rem:embodiment}, \ref{thm:embodied-jet}(1) says that the
memory line stores, once and simultaneously, the analytic family $\{(I^{z} f)(t)\}$ of
fractional integrals of the past --- and the jet of every order is a \emph{re-reading} of that
same stored family at shifted index, order $n$ living at Mellin depth
$\operatorname{Re} z > n$ of the strip. Nothing new is stored when one differentiates; the
embodiment of the signal already contains its jet, at every order below $z_* - 1$.
\end{remark}

% draft: Proposition 14.4. The statement ends at "order gamma"; the architecture
% comparison with the knot ladder, and the exactness of the identity beyond the strip
% condition (rational transforms continue meromorphically), are body text in the paper,
% moved out of the statement 2026-09-01 for verbatim sharing.
\begin{proposition}[the Gamma family: jets by subtraction along the shape column]\label{prop:gamma-jets}
\uses{prop:gamma-kernels,prop:gamma-cascade,thm:embodied-jet}
For the Gamma family write $u_m(\cdot, x)$ for the field of shape $m$ ($u_0 := f$). From
$s(1 + xs)^{-m} = x^{-1}\bigl[(1 + xs)^{-(m-1)} - (1 + xs)^{-m}\bigr]$,
\[
  x\,\partial_t\, u_m \;=\; u_{m-1} - u_m \;=\; -\,(\nabla u)_m,
  \qquad (\nabla u)_m := u_m - u_{m-1},
\]
and hence, iterating at fixed $x$,
\[
  x^n\,\partial_t^n\, u_\gamma \;=\; (-\nabla)^n u_\gamma
  \;=\; \sum_{j=0}^{n} (-1)^{n-j}\binom{n}{j}\, u_{\gamma - j}(t, x),
  \qquad 0 \le n \le \gamma .
\]
The column $\{u_m(\cdot, x)\}_{m=0}^{\gamma}$ is realized by $\gamma$ identical leaky
integrators in series with time constant $x$, $u_m = (1 + x\,\partial_t)^{-1} u_{m-1}$ --- the
gamma memory of \sscite{devries1992gamma} --- so the scale-normalized jet at scale $x$ is a
binomial-weighted \emph{subtraction of coexisting filter states}, with no temporal
differencing anywhere in the architecture, $\gamma$ states per scale, and jets available up
to order $\gamma$.
\statusT\quad\emph{(Transform algebra on \ref{prop:gamma-cascade}'s sections. The identity
is exact for all $n \le \gamma$, extending beyond \ref{thm:embodied-jet}'s strip condition
--- the transforms are rational and continue meromorphically --- consistently with
$z_* = \gamma$ marking the strip reading; the relaxation of (H) by meromorphic
continuation is among the open problems the paper records.)}
\end{proposition}

\begin{proof}
The transform identity is $(1+xs)^{-m}\bigl[(1+xs) - 1\bigr] = xs(1+xs)^{-m}$; induction in
$n$ at fixed $x$; the realization is \ref{prop:gamma-cascade} at $x_k = 0$,
where the section degenerates to the pure pole $(1 + x_{k+1}s)^{-1}$.
\end{proof}

% draft: Remark 14.5, as corrected in the paper (Rodrigues form; the draft's
% orthogonality phrasing was corrected on numerical verification, 2026-08).
\begin{remark}[receptive fields; Laguerre structure]\label{rem:laguerre}
\uses{prop:gamma-jets,prop:gamma-density}
By Rodrigues' formula the jet kernels are associated Laguerre: with $\phi^{(m)}_x$ the
Gamma density of shape $m$ at scale $x$ and $\phi_x = \phi^{(\gamma)}_x$,
\[
  x^n\,\partial_t^n\, \phi_x \;=\; \frac{n!\,\Gamma(\gamma - n)}{\Gamma(\gamma)}\,
  \phi^{(\gamma - n)}_x\; L_n^{(\gamma - n - 1)}\bigl(t/x\bigr),
  \qquad 0 \le n < \gamma,
\]
a Gamma density of shape $\gamma - n$ times an associated Laguerre polynomial --- the
shape dropping as the order rises, each derivative spending one of the kernel's poles.
And the shape column spans the order-$\gamma$ rational function space with the single
pole at $-1/x$: the classical Laguerre filter space at time constant $x$, exponentials
times polynomials of degree $< \gamma$, whose orthogonal basis is the Laguerre
functions. Both clauses are computations; the orthogonality belongs to the basis of the
space, not to the jet kernels themselves, whose polynomial parts all have degree
$\gamma - 1$. The comparison the clauses support is structural: the temporal
receptive-field family stands to the causal scale space
as the Hermite/derivative-of-Gaussian family stands to the Gaussian one --- a graded basis,
read off the embodied state by subtraction. The difference is that the grading here is
by the shape index rather than by differentiation order alone.
\end{remark}

% draft: Remark 14.6
\begin{remark}[the local corner: derivatives as curvatures of the memory trace]\label{rem:local-jet}
\uses{thm:embodied-jet,lem:moment-recursion,ex:bessel-quadratic,rem:gauge-freedom}
For the Bessel family, $z_* = \infty$ (\ref{lem:moment-recursion}(2)), so every order
is available; and in the parabolic gauge $\xi$ of \ref{rem:gauge-freedom},
\ref{thm:embodied-jet}(2) is the $n$-th power of the local operator of
\ref{ex:bessel-quadratic}:
\[
  \partial_t^n u \;=\; \Bigl(\tfrac12\,\partial_\xi^2 + \frac{\beta}{\xi}\,\partial_\xi\Bigr)^{n} u,
\]
the jet read from spatial stencils of order $2n$ on the instantaneous memory profile. In the
heat case $\beta = 0$ this is the statement that the observer reads temporal derivatives as
curvatures, and iterated curvatures, of its memory trace --- the embodiment reading already
implicit in \sscite{fagerstrom2005temporal}, Thm.~4, here extended from the local corner to
the whole class by the Mellin form.
\end{remark}


\appendix

\chapter{Memory kernels, Sonine conservation, and the scale-Cauchy problem: Theorem 3$'$}
% !TEX root = ../print.tex
% Draft \S9 --- Memory kernels, Sonine, Volterra
%
% Source: draft/hemigroup-causal-scale-space-kernels.md, lines 280-385.
% This file is chapter 9 --- Memory kernels, Sonine, Volterra of the blueprint; the shared counter makes every number here
% agree with the draft's own (see content.tex).
%
% Planned nodes: lem:memory-kernel, prop:scale-evolution, lem:potential-kernel, thm:sonine-conservation, cor:exact-inversion, prop:pair-regularity [A], prop:volterra; Example 9.9, Remarks 9.3/9.10/9.11
%
% No nodes yet. Each statement carries \label{kind:name}, \lean{...} with
% \leanok/\notready, \uses{...}, \notes{slug}, a \statusT or \statusA line, and a
% "% draft: <kind> N.M" provenance comment; an \statusA node must also declare what its
% citation carries (\textbf{Assignment.}).

% !TEX root = ../print.tex
% Draft \S10 --- The scale-Cauchy problem: Theorem 3'
%
% Source: draft/hemigroup-causal-scale-space-kernels.md, lines 386-463.
% Chapter 10 of the blueprint; the shared counter puts every number here on the draft's own.
%
% Theorem 3'. This chapter upgrades \ref{prop:scale-evolution} from a distributional identity
% to a classical Cauchy problem on an explicit operator core, and it is the first of the three
% primed theorems.
%
% \ref{thm:scale-cauchy} carries the longest proof in the article. Its two halves use quite
% different machinery and the formalisation should treat them separately: existence is Fubini
% plus the fundamental theorem of calculus for Bochner integrals (Mathlib has both), while
% uniqueness is a SCALAR linear ODE in the Laplace variable -- deliberately not a Hille--Yosida
% argument. That choice is what keeps the chapter inside the verified core: the only [A] node
% here is \ref{prop:fixed-scale-semigroup}, which is not used by \ref{thm:scale-cauchy} at all.
% It records the semigroup reading for its own sake.
%
% Notation warning carried from ledger A11: Phillips never writes -\varphi(-A), and he asserts
% D(B) contains D(A) rather than that D(A) is a core. The two citations in A11 are genuinely
% complementary, not redundant, and the node says so.

This chapter upgrades \ref{prop:scale-evolution} from a distributional identity to a classical
Cauchy problem in the scale variable, on an explicit operator core. Throughout we work in the
canonical gauge and write
\[
  \Xz := L^1(\RR_+),
\]
identified with the causal elements of $X$; by (A3) every $\Phi_{x,y}$ restricts to $\Xz$. On
$\Xz$ let $(T_r)_{r \ge 0}$ denote the \emph{delay semigroup},
$(T_r f)(t) = f(t - r)\mathbf{1}_{t \ge r}$, a strongly continuous semigroup of positive
isometries, and define the core
\[
  \core \;:=\; \bigl\{\, f \in \Xz :\ f \text{ absolutely continuous},\ f' \in \Xz,\ f(0) = 0 \,\bigr\}.
\]

% draft: Lemma 10.1
\begin{lemma}[delay derivative]\label{lem:delay-core}
\uses{lem:convolution-representation,def:cascade-family}
$\core$ is dense in $\Xz$ and invariant under every $T_r$ and every $\Phi_{x,y}$, and for
$f \in \core$,
\[
  \lim_{h \downarrow 0} \tfrac{1}{h}\bigl(T_h f - f\bigr) = -\,f' \quad \text{in } \Xz,
  \qquad
  \bigl\|T_r f - f\bigr\|_1 \le \min\bigl(2\|f\|_1,\ r\,\|f'\|_1\bigr).
\]
\statusT
\end{lemma}

\begin{proof}
Density is standard.

\emph{Invariance under $T_r$.} $T_r f$ is absolutely continuous, vanishes on $[0, r]$, and
$(T_r f)' = T_r f'$.

\emph{Invariance under $\Phi_{x,y} = \mu_{x,y} * \cdot$.} From $f = \mathbf{1}_{[0,\infty)} * f'$
we get $\mu * f = \mathbf{1}_{[0,\infty)} * (\mu * f')$, so $\mu * f$ is absolutely continuous
with derivative $\mu * f' \in \Xz$ and value $0$ at $0$.

\emph{The limit.} On $[h, \infty)$,
$\tfrac{1}{h}(f(t-h) - f(t)) + f'(t) = -\tfrac{1}{h}\int_{t-h}^{t}\bigl(f'(\rho) - f'(t)\bigr)d\rho$,
whose $L^1$-norm tends to $0$ by continuity of translation in $L^1$; on $[0, h)$ the
contribution is $\tfrac{1}{h}\int_0^h |f| + \int_0^h |f'| \le \sup_{[0,h]}|f| + o(1) \to
|f(0)| = 0$.

\emph{The estimate.} $T_r f - f = -\int_0^r T_\rho f'\,d\rho$ as a Bochner integral --- the
integrated form of the limit --- together with $\|T_\rho f'\|_1 = \|f'\|_1$ and the trivial
bound $\|T_rf - f\|_1 \le 2\|f\|_1$.
\end{proof}

% draft: Definition 10.2
\begin{definition}[per-scale generator; Phillips form]\label{def:phillips-generator}
\uses{lem:selfdecomposable-exponents,lem:delay-core}
Let $\nu_1$ be the Lebesgue--Stieltjes measure $-dk$ on $(0,\infty)$, with $k$ taken
right-continuous and $k(\infty) = 0$, so that $\nu_1\bigl((r,\infty)\bigr) = k(r)$; and for
$x > 0$ put $\nu_x := x^{-1}$ times the pushforward of $\nu_1$ under $r \mapsto xr$, i.e.\
$\nu_x\bigl((r,\infty)\bigr) = k(r/x)/x$. For $f \in \core$ define
\[
  \varphi_x(\partial_t)\, f \;:=\; b_0\, f' \;+\; \int_0^\infty \bigl(f - T_r f\bigr)\,\nu_x(dr)
  \;=\; b_0\, f' \;+\; \frac{1}{x}\int_0^\infty \bigl(f - T_{xr} f\bigr)\,\nu_1(dr).
\]
\statusT
\end{definition}

% draft: Lemma 10.3
\begin{lemma}[properties of the generator]\label{lem:generator-properties}
\uses{def:phillips-generator,lem:delay-core,lem:memory-kernel,prop:laplace-uniqueness}
For $f \in \core$ and $x > 0$:
\begin{enumerate}
\item The integral converges absolutely in $\Xz$, with
$\|\varphi_x(\partial_t) f\|_1 \le b_0 \|f'\|_1 + \int_0^\infty \min\bigl(2\|f\|_1,\, xr\,\|f'\|_1\bigr)\,x^{-1}\nu_1(dr) < \infty$.
\item $\Lap\bigl[\varphi_x(\partial_t) f\bigr](s) = \varphi_x(s)\,\hat f(s)$ for $s > 0$, where
$\varphi_x(s) = s F'(xs)$.
\item $\varphi_x(\partial_t)$ commutes with every $\Phi_{y,z}$ on $\core$.
\item $x \mapsto \varphi_x(\partial_t) f$ is continuous from $(0,\infty)$ to $\Xz$.
\item $\kappa^{(x)} * f$ is a.e.\ equal to an absolutely continuous function with
$\bigl(\kappa^{(x)} * f\bigr)' = \varphi_x(\partial_t) f$ a.e.; that is, the Phillips form
coincides on $\core$ with the memory-kernel operator of Chapter~9.
\end{enumerate}
\statusT
\end{lemma}

\begin{proof}
(1) Combine \ref{lem:delay-core} with $\int (1 \wedge r)\,\nu_x(dr) < \infty$, which follows
from $\int_0^1 k < \infty$ and $k(1) < \infty$ by integration by parts.

(2) Termwise: $\Lap[f'] = s\hat f$, using $f(0) = 0$, and $\Lap[f - T_rf] = (1 - e^{-sr})\hat f$.
Fubini is justified by (1), and $b_0 s + \int(1 - e^{-sr})\nu_x(dr) = \varphi_x(s)$ by
\ref{lem:memory-kernel} --- the two Bernstein representations of $\varphi_x$ agree.

(3) Convolutions commute; pull $\Phi_{y,z}$ through the Bochner integral.

(4) In the dilated form of \ref{def:phillips-generator} the integrand
$\tfrac1x(f - T_{xr}f)$ is, for $x$ in a compact subset of $(0,\infty)$, continuous in $x$ for
each $r$ by strong continuity of $T$, and dominated by $C\min(\|f\|_1,\, r\|f'\|_1)$, which is
$\nu_1$-integrable. Apply dominated convergence; the drift term is trivial.

(5) Both $\varphi_x(\partial_t)f \in L^1$ and $\kappa^{(x)} * f \in L^1_{\mathrm{loc}}$ are
causal with finite Laplace transforms. By (2) and \ref{lem:memory-kernel} the causal function
$V(t) := \int_0^t \bigl(\varphi_x(\partial_t)f\bigr)(\rho)\,d\rho - (\kappa^{(x)} * f)(t)$ has
transform $\varphi_x(s)\hat f(s)/s - F'(xs)\hat f(s) = 0$ for all $s > 0$, so $V = 0$ a.e.\ by
\ref{prop:laplace-uniqueness}. That is the claim; the value $(\kappa^{(x)}*f)(\zp) = 0$
follows from $\kappa^{(x)}([0,t]) \to b_0$ and $\sup_{[0,t]}|f| \to 0$.
\end{proof}

% draft: Theorem 10.4
\begin{theorem}[Theorem 3$'$: the scale-Cauchy problem]\label{thm:scale-cauchy}
\uses{thm:main-characterization,lem:delay-core,def:phillips-generator,lem:generator-properties,prop:scale-evolution,def:cascade-family,prop:laplace-uniqueness}
Let $(\Phi_{x,y})$ be as in \ref{thm:main-characterization}, in the canonical gauge, and let
$f \in \core$. Set $u(\cdot, x) := \Phi_{0,x} f$. Then:
\begin{enumerate}
\item $u(\cdot, x) \in \core$ for every $x \ge 0$;
\item $x \mapsto u(\cdot, x)$ belongs to
$C\bigl([0,\infty); \Xz\bigr) \cap C^1\bigl((0,\infty); \Xz\bigr)$ and satisfies
\[
  \partial_x u(\cdot, x) \;=\; -\,\varphi_x(\partial_t)\,u(\cdot, x) \quad (x > 0),
  \qquad u(\cdot, 0) = f,
\]
with $\partial_x$ a norm-derivative in $\Xz$;
\item $u$ is the \emph{unique} solution in the class of
$v \in C([0,\infty); \Xz) \cap C^1((0,\infty); \Xz)$ with $v(\cdot, x) \in \core$ for $x > 0$,
satisfying the equation for $x > 0$ and $v(\cdot, 0) = f$.
\end{enumerate}
\statusT\quad\emph{(the hemigroup analogue of \sscite{fagerstrom2005temporal}, Thm.~3, which
is recovered as the homogeneous case in \ref{rem:recovering-thm3}. Proved here, not cited.)}
\end{theorem}

\begin{proof}
(1) is the $\Phi$-invariance of $\core$ (\ref{lem:delay-core}), and continuity at every
$x \ge 0$ is (A7).

(2) Write $v(\cdot, x) := \varphi_x(\partial_t)\, u(\cdot, x)$. By
\ref{lem:generator-properties}(3), $v(\cdot, x) = \Phi_{0,x}\,\varphi_x(\partial_t) f$, so by
\ref{lem:generator-properties}(4), contractivity of $\Phi_{0,x}$, and (A7),
\[
  \|v(\cdot,x) - v(\cdot,x')\|_1
  \;\le\; \bigl\|\varphi_x(\partial_t)f - \varphi_{x'}(\partial_t)f\bigr\|_1
  + \bigl\|(\Phi_{0,x} - \Phi_{0,x'})\,\varphi_{x'}(\partial_t)f\bigr\|_1
  \;\longrightarrow\; 0,
\]
so $x \mapsto v(\cdot, x)$ is continuous from $(0,\infty)$ to $\Xz$.

By \ref{lem:generator-properties}(5) applied to $u(\cdot,x) \in \core$, each
$\kappa^{(x)} * u(\cdot, x)$ is absolutely continuous with derivative $v(\cdot, x)$, vanishing
at $0$; hence for $0 < a \le b$, Fubini gives
\[
  \int_a^b \bigl(\kappa^{(x)} * u(\cdot, x)\bigr)(t)\,dx
  \;=\; \int_0^t \int_a^b v(\rho, x)\,dx\,d\rho ,
\]
so the distributional $\partial_t$ in \ref{prop:scale-evolution}(2) evaluates to the Bochner
integral $\int_a^b v(\cdot, x)\,dx$:
\[
  u(\cdot, b) - u(\cdot, a) \;=\; -\int_a^b v(\cdot, x)\,dx .
\]
Since the integrand is continuous in $x$, the fundamental theorem of calculus for Bochner
integrals yields $u \in C^1((0,\infty); \Xz)$ with $\partial_x u = -v$, which is the equation.

(3) Let $v$ be any solution in the stated class and fix $s > 0$. Then $\eta(x) := \hat v(s, x)$
is continuous on $[0,\infty)$ and $C^1$ on $(0,\infty)$, with
\[
  \eta'(x) = \Lap[\partial_x v](s) = -\Lap[\varphi_x(\partial_t) v](s) = -\varphi_x(s)\,\eta(x)
\]
by \ref{lem:generator-properties}(2); differentiation and $\Lap$ exchange because
$|\hat g(s)| \le \|g\|_1$. Solving this scalar linear ODE on $[\varepsilon, x]$ and letting
$\varepsilon \downarrow 0$, using continuity at $0$, gives $\eta(x) = \hat f(s)\, e^{-F(xs)}$.
Thus $\hat v(s, x) = \hat u(s, x)$ for all $s > 0$, and $v = u$ by
\ref{prop:laplace-uniqueness}.
\end{proof}

% draft: Proposition 10.5
\begin{proposition}[fixed-scale semigroups; core property]\label{prop:fixed-scale-semigroup}
\uses{lem:potential-kernel,lem:delay-core,def:phillips-generator}
For fixed $x > 0$, the operator $-\varphi_x(\partial_t)$ with domain $\core$ is closable, and
its closure generates the strongly continuous contraction semigroup on $\Xz$ of convolutions
by the probability measures with Laplace transforms $e^{-\tau \varphi_x(s)}$, $\tau \ge 0$ ---
the semigroup subordinate to the delay semigroup via $\varphi_x$. Moreover $\core$ is a core
for that generator.
\statusA\quad\ledger{A11}\quad\emph{(Phillips' subordination theorem,
\sscite{phillips1952generation}, Thm.~4.3 with the hypotheses in its equation~(40); and the
core criterion, \sscite{engel2000one}, Prop.~II.1.7. Cited, not reproven here.
\textbf{Assignment.} \ledger{A11} carries two complementary halves and neither suffices alone.
Phillips gives the form of the subordinate generator; he asserts that its domain \emph{contains}
$\core$, not that $\core$ is a core for it, which is precisely the gap Engel--Nagel closes.
Two things the citation does not carry, both \textbf{[T]}: that $\core$ is dense and
semigroup-invariant, proved in \ref{lem:delay-core}, and that $\varphi_x$ is a Bernstein
function without killing term, proved in \ref{lem:potential-kernel}. A notational caution
recorded in the ledger: Phillips writes neither $-\varphi(-A)$ nor anything in that notation;
his displayed form becomes it only after setting the drift and killing parameters to zero.
Note finally that \ref{thm:scale-cauchy} does \emph{not} use this proposition --- the Cauchy
problem is solved by a scalar ODE, not by generation theory, so this node adds a reading
rather than a dependency.)}
\end{proposition}

\begin{proof}
$\varphi_x$ is a Bernstein function without killing term by \ref{lem:potential-kernel}, so
Phillips' subordination theorem applies to the delay semigroup, whose generator extends
$f \mapsto -f'$ on $\core$ (\ref{lem:delay-core}): the subordinate generator is the closure of
the Phillips form on the generator's domain. That $\core$ is a core follows from the cited
criterion, since $\core$ is dense and invariant under the subordinate semigroup, which
consists of convolutions (\ref{lem:delay-core}).
\end{proof}

% draft: Example 10.6
\begin{example}[the Gamma family: time-recursive form]\label{ex:gamma-recursive}
\uses{prop:gamma-family,def:phillips-generator}
For $F(s) = \gamma\log(1+s)$ we have $\nu_x(d\rho) = \gamma x^{-2} e^{-\rho/x}\,d\rho$, an
exponential memory, and $\varphi_x(s) = \gamma s/(1 + xs)$, so
\[
  \varphi_x(\partial_t)\,u \;=\; \gamma\,\partial_t\bigl(1 + x\,\partial_t\bigr)^{-1} u
  \;=\; \gamma\, w', \qquad \text{where } w \text{ solves } w + x\,w' = u,\ \ w(0) = 0 .
\]
The auxiliary variable $w$ is computed by one causal first-order ODE, so the scale step is
realized by a single recursive low-pass filter --- the strongest possible form of the
time-recursivity requirement of \sscite{fagerstrom2005temporal}, \S5: the observer embodies
its past in one first-order state per scale.
\end{example}

% draft: Remark 10.7
\begin{remark}[recovering Theorem 3 of 2005]\label{rem:recovering-thm3}
\uses{thm:scale-cauchy,prop:stable-family,rem:gauge-freedom}
For $F(s) = s^\alpha$ with $0 < \alpha < 1$ we have
$\nu_x\bigl((\rho,\infty)\bigr) = \alpha x^{\alpha-1}\rho^{-\alpha}/\Gamma(1-\alpha)$ and
$b_0 = 0$, so $\varphi_x(\partial_t) = \alpha x^{\alpha - 1} D^{\alpha}_{t,+}$ with
$D^{\alpha}_{t,+}$ the Riemann--Liouville/Marchaud derivative, and \ref{thm:scale-cauchy} reads
$\partial_x u = -\alpha x^{\alpha-1} D^{\alpha}_{t,+} u$. Under the reparametrization
$\tau := x^{\alpha}$ this is exactly
\[
  \partial_\tau u \;=\; -\,D^{\alpha}_{t,+}\, u,
\]
Theorem~3 of \sscite{fagerstrom2005temporal} verbatim --- identifying the scale parameter
$\tau$ of that paper as the $\alpha$-gauge of \ref{rem:gauge-freedom} applied to the canonical
scale coordinate.
\end{remark}

